% !TEX program = xelatex
% Методическое пособие «Введение в операционные системы и инструменты разработчика»
% Главный файл: собирает все главы в один PDF.
% openany: главы и части начинаются на следующей свободной странице,
% а не только на правой --- поэтому нет лишних пустых страниц.
\documentclass[12pt,a4paper,openany]{book}

\usepackage{style/course}

% ---- Титульный лист ----
\begin{document}
\frontmatter
\clearpage
\begin{titlepage}
  \begin{center}
    \vspace*{2.5cm}
    {\LARGE Новосибирский государственный университет\\[0.3em]
    Факультет информационных технологий\par}
    \vspace{3cm}
    {\Huge\bfseries
      Введение в операционные системы\\[0.4em]
      и инструменты разработчика\par}
    \vspace{1.5cm}
    {\Large Учебно-методическое пособие\par}
    \vspace{3cm}
    {\large М.В. Рутман\par}
    \vspace{2cm}
    {\Large 2026--2027\\[0.4em]
    1 курс ФИТ НГУ\par}
  \end{center}
\end{titlepage}

\setcounter{tocdepth}{2}
\tableofcontents
\clearpage

% ---- Основная часть ----
\mainmatter
\part{Введение и рабочая среда}
% Глава: 01-intro
% Статус: наполнено (первый прогон).
\chapter{Введение}
\label{ch:01-intro}

\section{О дисциплине}

Курс \emph{``Введение в операционные системы и инструменты разработчика''}
закладывает базовые навыки, которые потом используются на курсах
\emph{Объектно-ориентированное программирование},
\emph{Программная архитектура} и \emph{Операционные системы}.
В отличие от курса по программированию, здесь главный предмет --- не язык,
а \textbf{среда, в которой вы пишете и запускаете программы}:
командная оболочка, файловая система, процессы, компиляция, сборка,
взаимодействие процессов и виртуализация.

Цель дисциплины --- познакомить с основными инструментами разработки:
командными интерпретаторами и основами скриптинга, удалённым доступом,
системами управления версиями, сборщиками и отладчиками. По итогам курса
вы должны \emph{уметь} читать системную документацию, разбираться в том,
что делает система под капотом, и писать небольшие программы на C с их
сборкой, тестированием и отладкой.

Курс читается на первом курсе (оба семестра), без экзамена.
Контроль состоит из трёх частей:
\begin{itemize}
  \item \textbf{портфолио} --- накопленный набор выполненного: работа в
    терминале, собственные сценарии и программы, история коммитов;
  \item \textbf{групповой/итоговый проект} --- небольшое совместное изделие,
    собранный в единый репозиторий;
  \item \textbf{дифференцированный зачёт} --- один в каждом семестре,
    оценка по шкале \emph{неуд / уд / хорошо / отлично}.
\end{itemize}

\section{Рабочая среда курса: сервер лаборатории}

Основное рабочее место на курсе --- не ваш компьютер, а \textbf{сервер
лаборатории} \fpath{pglab.nsu.ru}. Курсу выделено отдельное
вычислительное устройство: на нём разворачивается по \textbf{одной
персональной виртуальной машине (ВМ) на каждого студента} --- изолированная
система на базе Ubuntu, доступная только по \textbf{SSH}. Развёртывание и
поддержание такого окружения --- самостоятельная инфраструктурная задача
университета; характеристики ВМ и настройка доступа --- в
\nameref{ch:02-linux-access}.

Вся практическая работа дисциплины выполняется именно на этом сервере:
работа в терминале и с файловой системой; написание, сборка, тестирование
и отладка программ на C (главы~\ref{ch:13-c-compilation}--\ref{ch:17-debugging});
контейнеры Docker (главы~\ref{ch:19-docker-basics}--\ref{ch:23-docker-resources});
Git и итоговый проект (\nameref{ch:25-git}).

Доступ к своей ВМ открыт в любое время: из кампуса --- напрямую, извне ---
через университетский \textbf{VPN} (см.\ \nameref{ch:02-linux-access}). По сути,
сервер лаборатории --- это «виртуальный компьютер» каждого студента на весь
семестр. Нагрузка на него существенна и устойчива: одновременно работающая
группа, у каждого --- своя ВМ, свои процессы, компиляции и контейнеры.

\begin{tipbox}
Локальный Linux на личном компьютере не требуется: вся среда --- на
\fpath{pglab.nsu.ru}. Подключаемся по SSH: \cmd{ssh <login>@pglab.nsu.ru}
(данные для входа выдаёт семинарист).
\end{tipbox}

\section{Как устроена методичка}

Методичка \textbf{тематическая}, а не ``по номерам семинаров''.
Одна тема может покрываться несколькими семинарами, а один семинар
--- несколькими темами. Главы собраны в пять частей:

\begin{center}
\begin{tabular}{ll}
\toprule
\textbf{Часть} & \textbf{Содержание} \\
\midrule
Часть 0  & Введение и рабочая среда \\
Часть I  & Операционная система и командная строка \\
Часть II & Разработка на C и сборка \\
Часть III& Взаимодействие процессов и виртуализация (Docker) \\
Часть IV & Среда разработчика (Git, CI, скриптинг) \\
\bottomrule
\end{tabular}
\end{center}

Каждая глава --- самостоятельный блок. Внутри главы есть
\emph{связь с работой ОС} (что происходит под капотом),
\emph{типичные ошибки}, \emph{ссылки на документацию}
(главным образом \cmd{man}) и \emph{задачи для самопроверки}.
Перекрёстные ссылки вида «см.\ главу~\ref{ch:03-shell}» ведут к связанным
темам. Состояние наполнения глав и их связь с лекциями и лабораторными
показывает \nameref{app:roadmap}.

\section{Как пользоваться методичкой}

\begin{itemize}
  \item Команды, флаги, переменные и термины написаны на английском
    (\cmd{ls -la}, \fpath{HOME}, \emph{container}) --- так их проще найти в
    документации и в интернете.
  \item Путь к файлу или директории оформлен как \fpath{/etc/passwd}.
  \item \emph{Задачи для самопроверки} в конце каждой главы --- это
    минимальный набор действий, которым вы должны владеть. Если что-то из них
    не удаётся выполнить --- перечитайте главу.
  \item Запоминать наизусть все флаги \textbf{не нужно}. Навык, который мы
    отрабатываем, --- быстро находить нужный параметр в документации
    (\nameref{ch:06-man}).
\end{itemize}

\begin{warnbox}
Если задача на семинаре не успевает быть решена, она автоматически
становится домашней. Методичка даёт материал и самопроверку, но выполнение
самостоятельно --- ваша работа.
\end{warnbox}

\section{Материалы и дальнейшее развитие}

Методичка растёт по мере прохождения курса: новые темы добавляются
по мере того, как появляются соответствующие семинары и лекции.
Главы, по которым материал ещё не развёрнут, помечены как
\emph{будет дополнено} и видны в \nameref{app:roadmap}.
Шпаргалки по командам собраны в \nameref{app:cheatsheets},
определения терминов --- в \nameref{app:glossary},
основная литература и ссылки --- в \nameref{app:references}.

\section{Задачи для самопроверки}
\begin{enumerate}
  \item Назовите, на какие последующие дисциплины опирается данный курс
    и какая у него основная цель.
  \item Какие три формы контроля составляют итоговую оценку и чем
    дифференцированный зачёт отличается от экзамена?
  \item Назовите пять частей методички и скажите, какой типичной теме
    посвящена каждая.
  \item Объясните, почему на курсе не требуется заучивать флаги команд
    наизусть и какой навык важнее.
\end{enumerate}

% Глава: 02-linux-access
% Статус: наполнено (первый прогон).
\chapter{Доступ к Linux-окружению}
\label{ch:02-linux-access}

В этой главе мы готовим рабочую среду. Основная модель курса ---
выданные вам виртуальные машины (ВМ) на университетском сервере,
к которым вы подключаетесь по SSH из собственного компьютера.
Локальная установка Linux (WSL, VirtualBox, нативная система)
нужна лишь как запасной вариант и рассмотрена в \ref{sec:local-linux}.

\section{Модель окружения}

Курсу выделяется вычислительный ресурс в виде сервера с хостнеймом
\fpath{pglab.nsu.ru}. На нём разворачиваются виртуальные машины
на базе Ubuntu; каждая ВМ --- отдельная изолированная система,
надёжно отделённая от других (об устройстве виртуализации и контейнеров
см.\ \nameref{ch:19-docker-basics}).

Характеристики выданного окружения:

\begin{itemize}
  \item \textbf{одна ВМ на студента} (по запросу --- больше);
  \item образ Ubuntu \textbf{без дополнительного ПО} --- пакеты ставите сами
    (\nameref{ch:08-packages});
  \item \textbf{20--40 ГБ} дискового пространства;
  \item срок жизни ВМ --- \textbf{до конца семестра};
  \item доступ к ВМ --- только по \textbf{SSH} (паролем или ключом),
    только из внутренней сети; извне --- через \textbf{VPN}
    (заявка на \fpath{support@nsu.ru}).
\end{itemize}

\begin{tipbox}
Работу выполняем в своей виртуальной машине, доступ к которой выдаётся на
сервере \fpath{pglab.nsu.ru}. Подключаемся по SSH:
\cmd{ssh <login>@pglab.nsu.ru} (данные для входа выдаёт семинарист).
Локальное окружение не требуется.
\end{tipbox}

На рис.\ \ref{fig:ssh-access} показана схема подключения.

\begin{figure}[htb]
  \centering
  \begin{tikzpicture}[node distance=1.2cm and 2.4cm, font=\small]
    \node[ext] (laptop) {Ваш компьютер\\[0.15em] Windows / macOS / Linux};
    \node[ext, below=of laptop] (vpn) {VPN-шлюз\\[0.15em] (из внешней сети)};
    \node[process, below=of vpn] (vm) {Ваша ВМ (Ubuntu)\\[0.15em]
        \fpath{pglab.nsu.ru}};
    \node[data, below=of vm] (fs) {Файловая система ВМ\\[0.15em]
        \fpath{/home/<login>}};
    \draw[dottedflow, ->] (laptop.south) -- node[right]{из внешней сети}
        (vpn.north);
    \draw[flow, ->] (vpn.south) -- node[right]{порт 22 (TCP)} (vm.north);
    \draw[flow, <->] (vm) -- (fs);
    \draw[flow, ->] (laptop.east) .. controls +(1.9,0) and +(1.9,0) ..
        node[right, pos=0.5]{SSH-сессия\ (терминал)} (vm.east);
  \end{tikzpicture}
  \caption{Подключение к рабочей ВМ. Из внутренней сети --- напрямую по SSH
    (порт 22); из внешней --- сначала через VPN. Сессия открывает терминал,
    работающий в файловой системе ВМ.}
  \label{fig:ssh-access}
\end{figure}

\section{Подключение по паролю}

SSH --- протокол безопасного удалённого доступа. Клиент
\cmd{ssh} уже установлен практически в любой системе
(проверьте \cmd{ssh -V}).

Проверка, что вы на месте:
\begin{minted}[bgcolor=gray!6]{bash}
ssh <login>@pglab.nsu.ru
\end{minted}

Если семинарист указал нестандартный порт, добавьте \flag{p}:
\begin{minted}[bgcolor=gray!6]{bash}
ssh -p <port> <login>@pglab.nsu.ru
# например: ssh -p 2222 student17@pglab.nsu.ru
\end{minted}

После установки соединения система запросит пароль.
Важная особенность: \textbf{при вводе пароль не отображается вообще}
(ни звёздочками, ни другими символами) --- это нормально.
Просто введите пароль и нажмите Enter.

Завершают сессию командой:
\begin{minted}[bgcolor=gray!6]{bash}
exit        # или нажмите Ctrl+D
\end{minted}

Важно понимать: ваши учебные файлы живут \textbf{на ВМ}, а не на
компьютере. Разрыв SSH-соединения \textbf{не удаляет} сохранённые файлы
--- при следующем входе они на месте (это следствие того, как
файловая система хранилась на диске ВМ, см.\ \nameref{ch:04-filesystem}).
Для входа нужен только доступ к сети.

\section{Подключение по ключу}

Пароль при каждом вводе неудобен и уязвим. Стандартная практика ---
авторизация по \textbf{ключевой паре} (эллиптическая/RSA-криптография).
Механизм асимметричный: у вас есть \emph{приватный ключ} (никому не
показывать) и \emph{публичный ключ} (разрешённый серверу).
При входе сервер шифрует вопрос так, что расшифровать его
может только приватный ключ владельца --- пароль не передаётся.

Генерируем пару на своём компьютере (команда \cmd{ssh-keygen};
разберитесь с ней через \cmd{man ssh-keygen}):
\begin{minted}[bgcolor=gray!6]{bash}
# по умолчанию ключи лежат в ~/.ssh/id_ed25519 и ~/.ssh/id_ed25519.pub
# ed25519: современный и короткий ключ; можно просто принять значения по умолчанию
ssh-keygen -t ed25519
\end{minted}

Копируем публичный ключ на ВМ в файл \fpath{~/.ssh/authorized_keys}.
Если доступ по паролю ещё работает, это делается одной командой:
\begin{minted}[bgcolor=gray!6]{bash}
ssh-copy-id <login>@pglab.nsu.ru
\end{minted}

Если \cmd{ssh-copy-id} недоступен (например, на Windows),
скопируйте публичный ключ \fpath{~/.ssh/id_ed25519.pub}
вручную. Затем допишите его содержимое в файл
\fpath{~/.ssh/authorized_keys} на ВМ
(в \nameref{ch:09-vim} разберём, как править файлы без
графического редактора).

После этого вход без пароля:
\begin{minted}[bgcolor=gray!6]{bash}
ssh <login>@pglab.nsu.ru
\end{minted}

\begin{warnbox}
Приватный ключ --- эквивалент пароля. Не копируйте
\fpath{id_ed25519} (без \fpath{.pub}) в облако и не коммитьте
в Git. Публичный ключ \fpath{id_ed25519.pub} передавать можно.
\end{warnbox}

\subsection{Файл конфигурации SSH}

Постоянно вводить логин, хост и порт не нужно --- опишем их
в конфигурации клиента \fpath{~/.ssh/config} (на вашем компьютере):
\begin{minted}[bgcolor=gray!6]{bash}
Host myvm
    HostName pglab.nsu.ru
    User  <login>
    Port  22
    IdentityFile ~/.ssh/id_ed25519
\end{minted}

Теперь достаточно:
\begin{minted}[bgcolor=gray!6]{bash}
ssh myvm
\end{minted}

\section{Доступ из внешней сети (VPN)}

ВМ доступна только из внутренней сети университета.
Работая из дома или из кафе, сначала подключаемся к VPN
(заявка на \fpath{support@nsu.ru}). После установления VPN-туннеля
подключение \cmd{ssh} работает как обычно --- сервер
становится достижимым по внутренней адресации.

\section{Локальный Linux как запасной вариант}
\label{sec:local-linux}

Если по какой-то причине вы не можете пользоваться сервером,
можно развернуть Linux локально. Основные варианты:
\begin{itemize}
  \item \textbf{WSL} (Windows Subsystem for Linux) --- подсистема,
    позволяющая запускать настоящий Linux из Windows без ВМ;
    устанавливается через «Приложение-программу» или
    \cmd{wsl --install -d Ubuntu} в PowerShell.
  \item \textbf{VirtualBox} --- классическая виртуальная машина;
    загружаем ISO Ubuntu и устанавливаем в ВМ.
  \item \textbf{Нативный Linux} (например, Ubuntu) на компьютере
    или на отдельном разделе.
\end{itemize}

Эти варианты эквивалентны для материала курса: всё, что описано
в остальных главах, выполняется одинаково. Но по умолчанию
мы предполагаем сервер.

\section{Типичные ошибки}

\begin{itemize}
  \item \emph{Connection timed out / No route to host} --- проверка
    сети, адреса и порта; при работе из дома --- установлен ли VPN.
  \item \emph{Permission denied} --- неверный логин или способ входа.
    Не перебирайте пароль многократно (счётчик попыток ведётся);
    покажите сообщение семинаристу.
  \item \emph{Ввёл пароль, но система говорит ``nope''} ---
    убедитесь, что не включён \fpath{CapsLock}, что вы вводите
    пароль для нужного пользователя.
  \item Работа «пропала» после перезагрузки компьютера ---
    вероятно, вы работали в локальной ВМ, а не на сервере;
    на сервере данные сохраняются.
\end{itemize}

\section{Задачи для самопроверки}
\begin{enumerate}
  \item Подключитесь к своей ВМ и выполните
    \cmd{hostname}, \cmd{id}, \cmd{df -h /} --- какие значения
    они выводят и что означает каждый?
  \item Сгенерируйте пару ключей:
    \cmd{ssh-keygen -t ed25519}. Затем настройте \fpath{~/.ssh/config}
    и убедитесь, что \cmd{ssh myvm} работает без запроса пароля.
  \item Объясните своими словами, чем асимметричная авторизация
    по ключу отличается от входа по паролю и почему приватный
    ключ нельзя передавать другим.
  \item Создайте файл \fpath{~/hello.txt} на ВМ, разрыв
    соединение, переподключитесь и убедитесь, что файл на месте.
    Почему он не исчез?
\end{enumerate}


\part{ОС и командная строка}
% Глава: 03-shell
% Статус: наполнено (первый прогон).
\chapter{Командная оболочка}
\label{ch:03-shell}

\section{Что такое оболочка}

\textbf{Оболочка} (shell) --- это программа-интерпретатор, которая читает
введённые вами команды, разбирает их и запускает соответствующие программы.
Наиболее распространённая оболочка Linux --- \cmd{bash}
(Bourne Again Shell). Оболочка --- это \emph{мост} между пользователем
и ядром операционной системы: она переводит ваш набор символов в вызовы
системных служб (\emph{syscalls}) и управляет жизненным циклом
запущенных программ.

\begin{sloppypar}
Все команды этой главы выполняются на вашей ВМ на сервере
\fpath{pglab.nsu.ru}.
\end{sloppypar}

\begin{figure}[htb]
  \centering
  \begin{tikzpicture}[node distance=1.6cm]
    \node[ext] (user) {Пользователь};
    \node[process, right=of user] (shell) {Оболочка\\[0.2em] bash};
    \node[process, right=of shell] (proc) {Запущенная\\[0.2em] программа};
    \node[process, below=of shell] (kernel) {Ядро Linux\\[0.2em]
        (syscalls)};
    \draw[flow, <->] (user) -- node[above, font=\small]{текст команд} (shell);
    \draw[flow, <->] (shell) -- node[right, font=\small]{fork/exec} (proc);
    \draw[flow, <->] (shell) -- node[left, font=\small]{open / read / write / exit} (kernel);
    \draw[flow, <->] (proc) -- (kernel);
  \end{tikzpicture}
  \caption{Оболочка парсит команды, создаёт подпроцессы
    (\emph{fork/exec}) и через ядро управляет ресурсами.}
  \label{fig:shell-kernel}
\end{figure}

Важно разделять \textbf{встроенные команды} (встроенные в саму оболочку:
\cmd{cd}, \cmd{echo}, \cmd{export}, \cmd{exit}) и \textbf{внешние
программы} (отдельные исполняемые файлы в \fpath{/usr/bin}, например
\cmd{ls}, \cmd{grep}). Встроенные команды работают быстрее и могут
менять состояние самой сессии (например, \cmd{cd} меняет каталог именно
\emph{текущей} оболочки).

\section{Сессии и подпроцессы}

Каждое окно терминала --- это отдельная \textbf{сессия} оболочки,
свой процесс. Запуск команды внутри оболочки порождает
\textbf{подпроцесс}. Подпроцесс наследует состояние родителя
(текущий каталог, переменные), но это \textbf{копия}: изменения
в подпроцессе не влияют на родителя.

\begin{minted}[bgcolor=gray!6]{bash}
cd ~/seminar1/source     # сменить каталог в текущей оболочке
bash                     # запустить вложенную оболочку: дочерний процесс
# внутри неё:
cd source                # сменится только каталог вложенной оболочки
touch tmp/x.txt          # файл создаётся в общей ФС и остаётся
exit                     # завершить вложенную оболочку
\end{minted}

После \cmd{exit} текущий каталог ``откатывается'' к значению
родительской оболочки, а созданный файл на месте --- файловая система
общая, а каталог --- состояние процесса.

\begin{tipbox}
Перейти в каталог \textbf{из скрипта} для всей текущей сессии
не получится: \cmd{cd} в скрипте выполняется в подпроцессе.
Поэтому для изменения окружения текущей сессии используют
\emph{подключение} файла (\cmd{source} или \cmd{.}),
см.\ раздел \ref{sec:bashrc}.
\end{tipbox}

\section{Переменные окружения}

\textbf{Переменная окружения} --- именованная строка, доступная
процессу. В оболочке переменная читается через \cmd{$NAME}:
\begin{minted}[bgcolor=gray!6]{bash}
echo $HOME        # домашний каталог
echo $USER        # текущий пользователь
echo $SHELL       # путь к оболочке по умолчанию
echo $PATH        # каталоги поиска программ
\end{minted}

\cmd{PATH} --- список каталогов, в которых оболочка ищет исполняемые
файлы при вводе команды. Именно благодаря \cmd{PATH} вы пишете
\cmd{ls}, а не \cmd{/usr/bin/ls}.

Создание переменной:
\begin{minted}[bgcolor=gray!6]{bash}
FOO=bar           # локальная переменная текущей оболочки
echo $FOO
\end{minted}

Локальная переменная \textbf{не видна} подпроцессам. Чтобы передать
её запускаемым программам, нужно \textbf{экспортировать}:
\begin{minted}[bgcolor=gray!6]{bash}
export FOO
# теперь дочерний процесс её видит:
bash -c 'echo $FOO'
\end{minted}

\cmd{export} переводит переменную в \emph{окружение} процесса,
которое наследуется дочерними. Список всех переменных окружения
выводит \cmd{env} или \cmd{printenv}; сравнить два окружения
удобно через \cmd{diff}.

\section{Файлы настройки оболочки}
\label{sec:bashrc}

Оболочка при запуске читает файлы настройки. Основные:

\begin{itemize}
  \item \fpath{/etc/profile} и \fpath{/etc/bash.bashrc} ---
    \emph{системные}, одинаковы для всех пользователей;
  \item \fpath{~/.profile} --- читается
    при \emph{входной} (login) сессии;
  \item \fpath{~/.bashrc} --- читается
    при \emph{интерактивной} сессии (открытие терминала).
\end{itemize}

На практике почти всё добавляют в \fpath{~/.bashrc}.
Например, чтобы переменная появлялась в каждой новой оболочке:
\begin{minted}[bgcolor=gray!6]{bash}
echo 'export COURSE="Operating Systems and Development Tools"' \
    >> ~/.bashrc
\end{minted}

Изменение вступит в силу в \textbf{новой} оболочке. Чтобы применить
сразу в текущей, подключите файл:
\begin{minted}[bgcolor=gray!6]{bash}
source ~/.bashrc     # или: . ~/.bashrc
\end{minted}

\cmd{source} выполняет команды \emph{в текущей} оболочке,
а не в подпроцессе --- именно поэтому переменная и остаётся.

\section{Типичные ошибки}
\begin{itemize}
  \item \emph{command not found} --- команда не найдена в \cmd{PATH}.
    Проверьте орфографию, установлен ли пакет; для полного пути
    используйте \cmd{which} / \cmd{type}.
  \item Изменение переменной ``не сохранилось'' в новый терминал ---
    вы создали локальную переменную и не экспортировали/не
    добавили в \fpath{.bashrc}.
  \item \emph{Permission denied} при запуске скрипта --- нет права
    на исполнение (см.\ \nameref{ch:05-permissions});
    можно запустить явно: \cmd{bash script.sh}.
  \item Скрипт ``меняет'' каталог/переменные, но после его запуска
    ничего не изменилось --- скрипт выполняется в подпроцессе;
    используйте \cmd{source}.
\end{itemize}

\section{Связь с остальным курсом}
Оболочка --- базис для всего остального: пакеты ставятся
командами (\nameref{ch:08-packages}), код пишется в редакторе
(\nameref{ch:09-vim}), сессии организуются в \nameref{ch:10-tmux},
а сценарии из команд собираются в
\emph{скрипты оболочки} (\nameref{ch:27-scripting-ide}).
Мониторинг процессов, запущенных оболочкой, ---
\nameref{ch:12-processes}.

\section{Задачи для самопроверки}
\begin{enumerate}
  \item Объясните разницу между встроенной командой и внешней
    программой; приведите по два примера.
  \item Создайте локальную переменную и покажите, что подпроцесс
    её не видит; затем экспортируйте и покажите, что видит.
  \item Разместите в \fpath{.bashrc} переменную, проверьте её в
    новой оболочке и объясните, почему \cmd{source} нужен,
    чтобы применить изменения сразу.
  \item Запустите вложенный \cmd{bash}, поменяйте в нём каталог
    и завершите его. Объясните, почему текущий каталог
    родительской оболочки не изменился, а созданный файл остался.
\end{enumerate}

% Глава: 04-filesystem
% Статус: наполнено (первый прогон).
\chapter{Файловая система}
\label{ch:04-filesystem}

\section{Единое дерево каталогов}

В Unix/Linux все ресурсы представляются как файлы, а файловая
система --- как \textbf{единое дерево} с одним корневым каталогом
\fpath{/}. В отличие от Windows (диски \fpath{C:}, \fpath{D:}),
здесь несколько физических дисков и разделов \emph{подключаются}
(\emph{mount}) в отдельные точки этого же дерева.

Ваш домашний каталог --- это \fpath{/home/<login>} на сервере
\fpath{pglab.nsu.ru}; всё, что вы создаёте, живёт на диске вашей ВМ.

Базовое устройство описывает стандарт \emph{FHS}
(Filesystem Hierarchy Standard). Ключевые каталоги:

\begin{center}
\begin{tabular}{ll}
\toprule
\textbf{Каталог} & \textbf{Назначение} \\
\midrule
\fpath{/etc}        & Конфигурация системных программ \\
\fpath{/home}       & Домашние каталоги пользователей \\
\fpath{/root}       & Домашний каталог \cmd{root} \\
\fpath{/usr}        & Основные программы и библиотеки системы \\
\fpath{/usr/bin}    & Исполняемые файлы (часть \fpath{PATH}) \\
\fpath{/var}        & Изменяемые данные: логи, кэши, почта \\
\fpath{/var/log}    & Файлы журналов системы и программ \\
\fpath{/tmp}        & Временные файлы (очищаются при перезагрузке) \\
\fpath{/dev}        & Файлы-устройства (диски, терминалы и т.п.) \\
\fpath{/proc}       & Виртуальная ФС с информацией о процессах и ядре \\
\fpath{/bin},\ \fpath{/sbin} & Критически важные команды (часто ссылки на \fpath{/usr}) \\
\bottomrule
\end{tabular}
\end{center}

На рис.\ \ref{fig:fs-tree} --- упрощённое представление дерева.

\begin{figure}[htb]
  \centering
  \begin{tikzpicture}[level distance=1.4cm, sibling distance=2.4cm]
    \node[process] {/}
      child {node[data] {/etc}}
      child {node[data] {/home}}
      child {node[data] {/usr}}
      child {node[data] {/var}}
      child {node[data] {/tmp}}
      child {node[data] {/dev}}
      child {node[data] {/proc}};
  \end{tikzpicture}
  \caption{Упрощённое дерево файловой системы Linux.}
  \label{fig:fs-tree}
\end{figure}

\section{Пути и навигация}

\begin{itemize}
  \item \fpath{/} --- корневой каталог;
  \item \fpath{~} --- домашний каталог текущего
    пользователя;
  \item \fpath{.} --- текущий каталог;
  \item \fpath{..} --- родительский каталог.
\end{itemize}

\textbf{Абсолютный путь} начинается с \fpath{/}
(\fpath{/home/alice/docs}); \textbf{относительный} --- от текущего
каталога (\fpath{docs/notes.txt}).

Основа навигации:
\begin{minted}[bgcolor=gray!6]{bash}
pwd                 # показать текущий каталог по абсолютному пути
ls                  # список содержимого
cd ~/work           # перейти в каталог
cd /                # перейти в корень
cd ..               # на уровень вверх
cd -                # вернуться в предыдущий каталог
ls -la ~/           # подробный список, включая скрытые файлы
\end{minted}

\cmd{ls} показывает \emph{видимое}; флаг \flag{l} ---
подробно (права, владелец, размер, время), \flag{a} --- включая
скрытые (имена, начинающиеся с \cmd{.}), \flag{h} --- размеры
в человекочитаемом виде, \flag{R} --- рекурсивно.
Все флаги всегда можно найти в \cmd{man ls}.

\section{Операции с файлами и каталогами}

\begin{center}
\begin{tabular}{lp{0.75\linewidth}}
\toprule
\textbf{Команда} & \textbf{Назначение} \\
\midrule
\cmd{mkdir}      & создать каталог (\flag{p} --- вместе с родителями) \\
\cmd{touch}      & создать пустой файл / обновить метки времени \\
\cmd{cp}         & копировать (\flag{r} --- рекурсивно, \flag{i} --- спросить, \flag{v} --- показывать) \\
\cmd{mv}         & переименовать или переместить \\
\cmd{rm}         & удалить файл (\flag{r} --- каталог с содержимым, \flag{i} --- спросить) \\
\cmd{rmdir}      & удалить \emph{пустой} каталог \\
\bottomrule
\end{tabular}
\end{center}

Типовые сценарии:
\begin{minted}[bgcolor=gray!6]{bash}
mkdir -p work/a/b/c     # создать цепочку каталогов за один шаг
touch work/a/b/c/one.txt work/a/b/c/two.txt
cp -rv work backup/     # показать, что копируется
mv one.txt renamed.txt  # переименовать
rm -i file.txt          # удалить с подтверждением
rm -rf work/a           # удалить всё дерево, осторожно!
mkdir empty && rmdir empty
\end{minted}

\begin{warnbox}
\cmd{rm -rf} на абсолютном пути не имеет ``корзины'' и не спрашивает.
Опечатка в пути (особенно в \fpath{/}) может удалить много данных.
Перед \cmd{rm -rf} перечитайте команду.
\end{warnbox}

\section{Метаданные, типы и размеры}

Файл в Linux --- это \emph{inode} (модель данных файловой системы,
см.\ далее) плюс имя. Команда \cmd{stat} показывает метаданные:
владельца, группу, права, размер, время изменения.
\cmd{file} определяет \textbf{тип} содержимого независимо от
расширения (по сигнатуре). \cmd{du} и \cmd{df} --- место:
\begin{minted}[bgcolor=gray!6]{bash}
stat config/service.conf       # метаданные
file image.bin                 # что это за файл на самом деле
du -sh ~/work                  # сколько места занимает каталог
df -h                          # свободное место на файловой системе
\end{minted}

\section{Поиск файлов}

Основной инструмент поиска --- \cmd{find}:
\begin{minted}[bgcolor=gray!6]{bash}
find . -name '*.conf'                 # все файлы с расширением .conf
find /home -type f -name 'user*'      # файлы, чьё имя начинается с user
find . -iname '*.log'                 # поиск без учёта регистра
find . -type d -name 'work'           # только каталоги
\end{minted}

\cmd{find} ищет по имени, типу, времени изменения и многим другим
атрибутам (\cmd{man find} --- одна из самых длинных справочных страниц).
Отдельная утилита \cmd{locate} ищет мгновенно по готовому индексу
имён, но индекс обновляется периодически.

\section{Как это устроено под капотом}

Файловая система хранит данные блоками на устройстве, а
\textbf{inode} --- это запись с атрибутами файла: права, владелец,
размер, указатели на блоки данных. Имя файла живёт в каталоге
(каталог --- это тоже файл, список «имя $\to$ inode»).
Это объясняет две вещи:
\begin{itemize}
  \item жёсткая связь (\emph{hard link}) --- ещё одно имя
    на тот же inode; удаление имени не удаляет данные, пока
    есть другие имена;
  \item \cmd{mv} внутри одной ФС очень дёшев --- перемещается
    только имя (запись в каталоге), а данные на месте.
\end{itemize}

Механизм inode объясняет и то, почему данные на вашей ВМ
(см.\ \nameref{ch:02-linux-access}) сохраняются после
разрыва SSH: они записаны в блочное хранилище ВМ,
а сессия --- лишь процесс, работающий с этим хранилищем.

\section{Типичные ошибки}
\begin{itemize}
  \item \emph{No such file or directory} --- опечатка в пути или
    файл не в текущем каталоге (забудьте \fpath{./}).
  \item Скрытые файлы ``не видно'' в \cmd{ls} --- нужен \flag{a}.
  \item \cmd{rmdir} не удаляет каталог --- он не пуст; используйте
    \cmd{rm -r} (после проверки содержимого).
  \item \cmd{mv a b}, где \fpath{b} --- каталог, перемещает
    \fpath{a} \emph{внутрь} \fpath{b}, а не переименовывает.
  \item Ищите файл, а \cmd{locate} не находит --- индекс устарел;
    используйте \cmd{find}.
\end{itemize}

\section{Связь с остальным курсом}
Права и владельцы, описанные в метаданных, --- тема
\nameref{ch:05-permissions}. Поиск и текстовые утилиты для анализа
файлов --- \nameref{ch:07-text-pipelines}. Монтирование точек
(откуда берётся дерево из нескольких дисков) и контейнеры,
изолирующие файловую систему, --- \nameref{ch:19-docker-basics}
и \nameref{ch:23-docker-resources}.

\section{Задачи для самопроверки}
\begin{enumerate}
  \item Объясните разницу между абсолютным и относительным путём;
    чем отличается \fpath{~} от \fpath{.}.
  \item Одной командой создайте каталог с тремя вложенными уровнями
    и пустой файл в самом глубоком.
  \item Определите тип произвольного файла, не смотря на расширение.
    Как вы это сделали?
  \item Найдите все файлы с расширением \fpath{.txt} в домашнем
    каталоге и подсчитайте их количество.
  \item Объясните на примерах, почему \cmd{mv} внутри одной
    файловой системы не копирует данные.
\end{enumerate}


% Глава: 05-permissions
% Статус: наполнено (первый прогон).
\chapter{Права доступа и пользователи}
\label{ch:05-permissions}

\section{Пользователи и группы}

Linux --- многопользовательская система: каждый процесс работает от
имени \textbf{пользователя} и принадлежит \textbf{группе}. Каждому
пользователю присвоено числовое идентификатор (UID), каждой группе ---
GID. Текстовые имена хранятся в \fpath{/etc/passwd} (пользователи)
и \fpath{/etc/group} (группы).

На сервере \fpath{pglab.nsu.ru} вы --- обычный пользователь своей ВМ;
привилегии суперпользователя получают только администраторы.

Свой аккаунт и группы:
\begin{minted}[bgcolor=gray!6]{bash}
id                     # uid, gid и все группы текущего пользователя
whoami                 # имя текущего пользователя
groups                 # группы
getent passwd          # список пользователей, включая LDAP
\end{minted}

Пользователь с UID~0 --- \cmd{root} (администратор). Обычные
пользователи не могут менять системные файлы, пока не получат
временные привилегии.

\section{Права на файл}

У каждого файла (и каталога) есть \textbf{владелец},
\textbf{группа} и \textbf{права}. Права задаются для трёх субъектов:
\textbf{владелец} (owner), \textbf{группа} (group),
\textbf{остальные} (others). Для каждого субъекта возможны
три операции: \textbf{чтение} (r), \textbf{запись} (w),
\textbf{исполнение} (x).

Разбор записи \cmd{ls -l}:
\begin{minted}[bgcolor=gray!6]{bash}
-rw-r--r-- 1 alice alice 1234 Jan 1 12:00 notes.txt
drwxr-xr-x 2 alice devgrp    4096 Jan 1 12:00 project
\end{minted}

Первый символ --- тип файла: \cmd{-} обычный, \cmd{d} каталог,
\cmd{l} символическая ссылка. Затем три группы по три буквы:
владелец, группа, остальные. В примере владелец читает и пишет
\emph{и исполняет}, группа и остальные --- только читают.

Значение \textbf{исполнения} зависит от типа:
для файла --- можно запустить как программу; для каталога ---
можно войти внутрь (\cmd{cd}) и видеть содержимое.
Без \cmd{x} на каталоге нельзя перейти в него, даже если есть \cmd{r}.

\section{Числовая запись и \cmd{chmod}}

Права удобно задавать числом: для каждого субъекта
\cmd{r=4}, \cmd{w=2}, \cmd{x=1}, сумма --- одна цифра.
Например \cmd{7} (rwX), \cmd{5} (r-x), \cmd{6} (rw-), \cmd{4} (r--).
\begin{minted}[bgcolor=gray!6]{bash}
chmod 755 script.sh     # rwxr-xr-x
chmod 644 notes.txt     # rw-r--r--
chmod +x script.sh      # дать право выполнения всем
chmod u+w file.txt      # дать владельцу право записи
chmod a-x file.txt      # убрать право выполнения у всех
\end{minted}

\textbf{umask} --- маска, из которой система вычитает права
при создании файлов: это объясняет, почему новый файл
получает \cmd{644}, а новый каталог --- \cmd{755}.
\cmd{umask} без аргументов показывает текущую маску.

\section{Смена владельца и группы}

Собственником файлов можно стать только от имени \cmd{root}
(обычно через \cmd{sudo}, см.\ ниже).
\begin{minted}[bgcolor=gray!6]{bash}
sudo chown alice:devgrp file.txt   # сменить владельца и группу
sudo chown alice file.txt          # только владелец
chgrp devgrp file.txt              # группа, если вы в этой группе
\end{minted}

\section{\cmd{sudo}: привилегии администратора}

\cmd{sudo} выполняет команду \textbf{от имени другого пользователя}
(по умолчанию \cmd{root}) после проверки. Это основной способ
делать системные действия (установка пакетов, смена владельца):
\begin{minted}[bgcolor=gray!6]{bash}
sudo apt update            # выполнить команду от root
sudo -u postgres psql      # выполнить команду от пользователя postgres
sudo su -                  # получить оболочку root
\end{minted}

После первой проверки пароля в течение нескольких минут
\cmd{sudo} не спрашивает его повторно. Пользователь
попал в список, кому разрешён \cmd{sudo}, --- в
\fpath{/etc/sudoers} (редактируется через \cmd{visudo}).

\begin{warnbox}
Внутри сессии под \cmd{root} (после \cmd{sudo su -})
легко ошибиться с путём при \cmd{rm}. Для одноразовых
системных команд лучше \cmd{sudo <команда>},
чем входить в root-оболочку.
\end{warnbox}

\section{Частые специальные права (кратко)}

Существуют расширенные биты:
\textbf{setuid} и \textbf{setgid} (файл выполняется
с правами владельца/группы файла) и \textbf{sticky bit}
на каталоге (удалять файлы в нём может только владелец ---
так устроено \fpath{/tmp}). Они видны как \cmd{s}, \cmd{S},
\cmd{t} в записи \cmd{ls -l}. Для курса достаточно знать их
наличие и где они встречаются (например, \cmd{su}, \cmd{passwd}
имеют setuid).

\section{Типичные ошибки}
\begin{itemize}
  \item \emph{Permission denied} при запуске скрипта --- нет \cmd{x};
    дайте \cmd{chmod +x}.
  \item Нельзя открыть файл в каталоге --- нет \cmd{r} на каталоге
    \emph{или} нет \cmd{x} на нём (нужны оба).
  \item \cmd{chmod} не сработал --- вы не владелец и не \cmd{root}.
  \item Скрипт с \cmd{\#!/bin/bash} не запускается --- нет \cmd{x}
    или неверная строка интерпретатора в первой строке.
  \item Не работает \cmd{sudo} --- пользователь не в группе,
    которой разрешён sudo (проверьте с преподавателем).
\end{itemize}

\section{Связь с остальным курсом}
Права на файлы --- основа безопасности. Монтирование в Docker
(\nameref{ch:21-docker-data}) работает по той же модели:
что контейнеру доступно, определяется правами и точками монтирования.
Пакеты ставятся с правами \cmd{root}
(\nameref{ch:08-packages}), а работа под другим пользователем
(\cmd{sudo -u}) часто нужна при запуске сервисов.

\section{Задачи для самопроверки}
\begin{enumerate}
  \item Расшифруйте запись \fpath{drwxr-x---}: тип файла, права
    владельца, группы, остальных.
  \item Числовой записью задайте права «владелец читает/пишет,
    остальные --- только чтение». Какую команду вы выполнили?
  \item Создайте каталог, в который вам нельзя войти, и объясните
    по каким правам это произошло.
  \item Объясните разницу между \cmd{chmod u+w} и \cmd{chmod 666}.
  \item Как выполнить команду один раз от имени другого
    пользователя, не заходя в его оболочку?
\end{enumerate}

% Глава: 06-man
% Статус: наполнено (первый прогон).
\chapter{Документация: man, info, help}
\label{ch:06-man}

\section{Зачем это важно}

В Linux основная документация по программам лежит \emph{на самой
системе}: команды сопровождаются справочными страницами, которые
открываются утилитой \cmd{man}. Научиться читать \cmd{man} ---
навык важнее, чем заучивать команды: справочник позволяет
разобраться в любой утилите, даже о которой вы не слышали.

Все справочные страницы доступны прямо на вашей ВМ на
\fpath{pglab.nsu.ru}: \cmd{man} читает документацию, установленную
на сервере.

\begin{tipbox}
Не пытайтесь запомнить все флаги. Запомните \emph{паттерн}:
«откройте \cmd{man}, найдите нужный флаг, прочитайте его
описание». Это покрывает 90\% бытовых задач.
\end{tipbox}

\section{Базовое использование}

\begin{minted}[bgcolor=gray!6]{bash}
man ls              # открыть справочную страницу
man -k "regular expression"   # поиск по кратким описаниям (apropos)
man -S 2 ls         # конкретный раздел
man bash            # документация по оболочке
\end{minted}

Внутри \cmd{man}:
\begin{itemize}
  \item \cmd{q} --- выход;
  \item \cmd{/pattern} --- поиск, \cmd{n} / \cmd{N} --- следующее/
    предыдущее совпадение;
  \item \cmd{g} --- в начало, \cmd{G} --- в конец;
  \item пробел / \cmd{b} --- страница вперёд / назад.
\end{itemize}

Этот интерфейс --- стандартная \textbf{pager} (\cmd{less});
он же используется при чтении больших файлов
(см.\ \nameref{ch:07-text-pipelines}).

\section{Структура справочной страницы}

Страница разбита на разделы. Типовой порядок:
\begin{center}
\begin{tabular}{lp{0.6\linewidth}}
\toprule
\textbf{Раздел} & \textbf{Содержимое} \\
\midrule
NAME       & одно-два предложения: что делает программа \\
SYNOPSIS   & краткий синтаксис, часто с основными флагами \\
DESCRIPTION& подробное описание работы \\
OPTIONS    & все флаги и их смысл (ключевой раздел!) \\
FILES      & файлы, с которыми работает программа \\
EXAMPLES   & готовые примеры использования \\
SEE ALSO   & связанные страницы и утилиты \\
\bottomrule
\end{tabular}
\end{center}

Самый полезный раздел для быстрой работы --- \cmd{OPTIONS}:
он содержит все флаги в алфавитном порядке.

\section{Секции man}

Справочные страницы разделены на \textbf{секции} с цифрами
(от 1 до 9):

\begin{center}
\begin{tabular}{cl}
\toprule
\textbf{Секция} & \textbf{Содержимое} \\
\midrule
1 & Пользовательские команды (\cmd{ls}, \cmd{grep}, \cmd{git}) \\
2 & Системные вызовы ядра (\cmd{open}, \cmd{read}, \cmd{fork}) \\
3 & Библиотечные функции (C: \cmd{malloc}, \cmd{printf}) \\
4 & Файлы устройств (\fpath{/dev/console}) \\
5 & Форматы файлов и конфиги (\fpath{/etc/passwd}) \\
8 & Административные команды (\cmd{mount}, \cmd{useradd}) \\
\bottomrule
\end{tabular}
\end{center}

Если одна команда есть в нескольких секциях, в заголовке
пишут \fpath{ls (1)}, \fpath{ls (8)} и т.п. По умолчанию
открывается первая доступная.

\section{info и встроенный \cmd{--help}}

Помимо \cmd{man} существуют другие источники:
\begin{itemize}
  \item \cmd{--help} (или \cmd{-h}) у большинства программ
    --- краткая справка в терминале, без пагина;
  \item \cmd{info <prog>} --- гиперссылочная документация
    (GNU-инструменты: \cmd{make}, \cmd{gcc}, \cmd{git});
  \item \cmd{apropos <term>} --- поиск по кратким описаниям
    (синоним \cmd{man -k});
  \item \cmd{whatis <prog>} --- только раздел NAME.
\end{itemize}

Правило выбора: для быстрой проверки флага ---
\cmd{--help}; для систематического изучения ---
\cmd{man} или \cmd{info}.

\section{Типичные ошибки}
\begin{itemize}
  \item \emph{No manual entry for X} --- программа не
    установлена или у неё нет страницы; проверьте
    \cmd{which X}.
  \item Не выходит из \cmd{man} --- нажмите \cmd{q}.
  \item Искать весь список флагов в интернете, когда
    \cmd{man} стоит на вашей машине.
  \item Пытаться читать \cmd{man} без понимания структуры:
    сначала \cmd{NAME}, потом \cmd{SYNOPSIS}, и только
    затем \cmd{OPTIONS}.
\end{itemize}

\section{Связь с остальным курсом}

\cmd{man} --- универсальный инструмент, который пригодится
везде: в работе с пакетами
(\nameref{ch:08-packages}), C-компилятором
(\nameref{ch:13-c-compilation}), \cmd{git}
(\nameref{ch:25-git}), Docker (\nameref{ch:19-docker-basics}).
Секция~2 --- системные вызовы, о которых подробнее в главе
о процессах (\nameref{ch:12-processes}).

\section{Задачи для самопроверки}
\begin{enumerate}
  \item Найдите в \cmd{man} флаг \cmd{ls}, который показывает
    размеры в человекочитаемом виде. Как он называется?
  \item Объясните назначение секций 1, 2, 5, 8 и приведите
    по одному примеру каждой.
  \item Как быстро узнать, что делает незнакомая команда,
    не читая всю страницу?
  \item Чем \cmd{--help} отличается от \cmd{man} по
    объёму и структуре?
  \item Найдите в документации \cmd{grep} флаг, который
    выводит номер строки.
\end{enumerate}

% Глава: 07-text-pipelines
% Статус: наполнено (первый прогон).
\chapter{Стандартные потоки, перенаправление, конвейеры}
\label{ch:07-text-pipelines}

\section{Стандартные потоки}

Каждый запускаемый процесс по умолчанию получает три
\textbf{стандартных потока}:

\begin{center}
\begin{tabular}{clp{0.4\linewidth}}
\toprule
\textbf{FD} & \textbf{Поток} & \textbf{Назначение} \\
\midrule
0 & stdin  & вход: то, что программа читает \\
1 & stdout & выход: результат работы программы \\
2 & stderr & ошибки и диагностика \\
\bottomrule
\end{tabular}
\end{center}

По умолчанию все три подключены к терминалу. Ключевой принцип
Unix: программы пишут результат в \cmd{stdout}, а сообщения
об ошибках --- в \cmd{stderr}. Разделение позволяет
независимо перенаправлять данные и диагностику.

\section{Перенаправление вывода}

\begin{minted}[bgcolor=gray!6]{bash}
ls > out.txt        # stdout в файл: создаёт или очищает
ls >> out.txt       # stdout добавляется в файл
ls 2> err.txt       # stderr в файл
ls &> all.txt       # stdout и stderr в один файл (bash)
ls 2>&1 > out.txt   # stderr туда, где сейчас stdout, потом stdout в файл
ls > out.txt 2> err.txt    # явно разделить
tee out.txt < big.txt       # показать на экране и в файле
\end{minted}

Запомнить два оператора: \cmd{>} --- заменить/создать файл,
\cmd{>>} --- дописать.

\begin{warnbox}
\cmd{>} \textbf{обрезает} файл до нуля перед записью. Если
файл уже содержал важные данные и вы хотели дописать,
используйте \cmd{>>}.
\end{warnbox}

\section{Конвейеры (pipes)}

Оператор \cmd{|} передаёт \cmd{stdout} одной команды в
\cmd{stdin} другой. Это основной способ строить сложные
преобразования из простых:
\begin{minted}[bgcolor=gray!6]{bash}
cat access.log | grep "GET" | wc -l
df -h | grep " /$"
ps aux | sort -k3 -nr | head
\end{minted}

Правило: каждое звено конвейера делает \textbf{одну} операцию.
Если решение требует трёх отдельных команд с временными
файлами --- скорее всего, его можно выразить конвейером.

\section{Текстовые утилиты}

\begin{center}
\begin{tabular}{lp{0.6\linewidth}}
\toprule
\textbf{Утилита} & \textbf{Назначение} \\
\midrule
\cmd{cat}   & вывести файл(ы) в stdout; склеить файлы \\
\cmd{head}  & первые строки/байты (\flag{n} 5 --- 5 строк) \\
\cmd{tail}  & последние строки; \flag{f} --- следить за обновлением \\
\cmd{wc}    & подсчитать строки/слова/байты (\flag{l}, \flag{w}, \flag{c}) \\
\cmd{sort}  & отсортировать строки (\flag{r} --- обратно, \flag{n} --- число, \flag{k} --- ключ) \\
\cmd{uniq}  & убрать подряд идущие дубли (работает после sort) \\
\cmd{cut}   & вырезать поле по разделителю (\flag{-d} '\textbackslash|' \flag{-f} 2) \\
\cmd{grep}  & поиск по регулярному выражению \\
\cmd{tr}    & побайтовая замена/удаление символов \\
\cmd{less}  & постраничный просмотр больших файлов \\
\bottomrule
\end{tabular}
\end{center}

\cmd{grep} --- самая частая. Основные флаги:
\begin{minted}[bgcolor=gray!6]{bash}
grep ERROR log.txt         # строки, содержащие ERROR
grep -i error log.txt      # без учёта регистра
grep -c ERROR log.txt      # сколько таких строк
grep -E "err|warn" log.txt # расширенные регулярные выражения
grep -v DEBUG log.txt      # строки, НЕ содержащие DEBUG
grep -n "user" config      # с номерами строк
\end{minted}

\section{Типовой сценарий анализа}

\begin{sloppypar}
На практике такой лог-файл обычно лежит на вашей ВМ на сервере
\fpath{pglab.nsu.ru}. Допустим, в \fpath{events.log} строки вида
\cmd{timestamp|level|source|file|status}.
\end{sloppypar}
Частые вопросы:
\begin{minted}[bgcolor=gray!6]{bash}
wc -l < events.log                       # общее число строк
grep -c ERROR events.log                 # число ошибок
cut -d'|' -f2 events.log | sort | uniq -c # количество каждого уровня
cut -d'|' -f3 events.log | sort | uniq -c | sort -nr # по источникам
grep ERROR events.log > errors.log        # выписать ошибки
\end{minted}

Ключевой приём: \cmd{sort} перед \cmd{uniq} ---
\cmd{uniq} убирает только \emph{соседние} дубли, поэтому
сначала строки должны быть отсортированы.

\section{Постраничный просмотр: \cmd{less}}

Для больших файлов \cmd{cat} неудобен: вывод проматывает
терминал. \cmd{less} --- стандартный pager:
\begin{minted}[bgcolor=gray!6]{bash}
less big.log
# внутри less:
# /pattern  n  N  - поиск вперёд/назад
# g  G            - начало / конец
# Space / b       - страница вперёд / назад
# q               - выход
\end{minted}

\cmd{less} и \cmd{more} похожи, но \cmd{less} позволяет
прокручивать назад и искать.

\section{Типичные ошибки}
\begin{itemize}
  \item Конвейер «ничего не возвращает», хотя отдельные
    команды работают --- проверьте, не фильтрует ли
    \cmd{grep} всё (нет совпадений);
  \item \cmd{uniq} не убирает все дубли --- не отсортированы
    строки перед \cmd{uniq};
  \item \cmd{>} случайно перезаписал важный файл;
  \item Забыли \cmd{2>} и диагностика ушла в терминал,
    а не в файл;
  \item \cmd{tail} без \flag{f} показывает конец один раз
    и завершается; для наблюдения нужен \flag{f}.
\end{itemize}

\section{Связь с остальным курсом}

Конвейеры --- это основа автоматизации: из них собираются
скрипты оболочки (\nameref{ch:27-scripting-ide}),
обрабатываются логи приложений, выводятся результаты
тестов (\nameref{ch:15-testing-optimization}), логируются
запуски контейнеров (\nameref{ch:19-docker-basics}).

\section{Задачи для самопроверки}
\begin{enumerate}
  \item Объясните разницу между \cmd{>} и \cmd{>>}. Что
    произойдёт, если вы случайно использовали первое?
  \item Что такое \cmd{stderr} и зачем его разделять от
    \cmd{stdout}?
  \item Напишите конвейер, который считает количество
    строк в файле, содержащих слово \fpath{error}
    (без учёта регистра).
  \item Почему \cmd{sort} обычно ставят перед \cmd{uniq}?
  \item Как посмотреть последние 20 строк журнала и
    продолжать наблюдать за новыми?
\end{enumerate}

% Глава: 08-packages
% Статус: наполнено (первый прогон).
\chapter{Пакетная система: dpkg и apt}
\label{ch:08-packages}

\section{Зачем нужна пакетная система}

Устанавливать программы «вручную» (скачать архив, распаковать,
поставить зависимости) неудобно и ненадёжно.
\textbf{Пакетная система} хранит программы в стандартизированном
виде и отслеживает: что установлено, какие файлы к чему
относятся, какие зависимости есть. В Debian/Ubuntu
основной формат --- \textbf{\fpath{.deb}}.

Устанавливаем пакеты в свою ВМ на сервере \fpath{pglab.nsu.ru}:
образ Ubuntu идёт «голый», и каждый добирает нужное сам.

\begin{figure}[htb]
  \centering
  \begin{tikzpicture}[node distance=1.4cm]
    \node[process] (repo) {Депозитарий\\[0.2em] (apt)};
    \node[data, below=of repo] (cache) {Кэш пакетов\\[0.2em] /var/cache/apt};
    \node[process, below=of cache] (dpkg) {dpkg\\[0.2em] (распаковка .deb)};
    \node[storage, below=of dpkg] (fs) {Файлы системы\\[0.2em] /usr, /etc};
    \draw[flow, <->] (repo) -- node[right]{установка} (cache);
    \draw[flow, <->] (cache) -- (dpkg);
    \draw[flow, <->] (dpkg) -- (fs);
  \end{tikzpicture}
  \caption{Связь между apt (репозиторий) и dpkg (развёртка).}
  \label{fig:pkg}
\end{figure}

\section{dpkg и apt: чем отличаются}

\begin{center}
\begin{tabular}{lp{0.4\linewidth}lp{0.4\linewidth}}
\toprule
\textbf{dpkg} & работает с конкретным файлом \fpath{.deb} и файлами
системы; зависимости сам не разрешает \\
\textbf{apt} & управляет набором пакетов с учётом зависимостей и репозиториев \\
\midrule
\multicolumn{2}{l}{На практике: \cmd{apt} --- интерфейс для пользователя,} \\
\multicolumn{2}{l}{\cmd{dpkg} --- низкоуровневый исполнитель.} \\
\bottomrule
\end{tabular}
\end{center}

Обычно работают через \cmd{apt} (или \cmd{apt-get}); \cmd{dpkg}
используют для точечных задач (установка локального \fpath{.deb},
просмотр файлов пакета).

\section{Основные операции}

\begin{minted}[bgcolor=gray!6]{bash}
sudo apt update              # обновить список доступных пакетов
sudo apt upgrade             # обновить установленные пакеты
sudo apt install <pkg>       # установить пакет
sudo apt remove <pkg>        # удалить, файлы конфигурации остаются
sudo apt purge <pkg>         # удалить вместе с конфигурацией
sudo apt autoremove          # удалить ненужные зависимости
apt search <term>            # поиск по каталогу
apt show <pkg>               # сведения о пакете
apt list --installed         # список установленных пакетов
\end{minted}

\textbf{update} и \textbf{upgrade} --- разные вещи:
\cmd{update} перечитывает списки доступных версий из
репозиториев; \cmd{upgrade} ставит новые версии уже
установленного.

\section{Работа с локальными \fpath{.deb}}

Если вам дали файл \fpath{package.deb} --- его можно установить
двумя способами:
\begin{minted}[bgcolor=gray!6]{bash}
sudo dpkg -i package.deb          # напрямую
sudo apt install ./package.deb    # apt разрешает зависимости
\end{minted}

Первый способ быстрее, но если у пакета есть
зависимости, которых нет в системе, \cmd{dpkg}
остановится. Тогда \cmd{apt install ./package.deb}
доустановит всё из репозиториев.

\section{Инспекция пакетов без установки}

\begin{minted}[bgcolor=gray!6]{bash}
dpkg-deb --info package.deb      # имя, версия, описание
dpkg-deb --contents package.deb  # список файлов внутри
dpkg-deb --show --showformat='...\\n' package.deb  # свои поля
\end{minted}

После установки:
\begin{minted}[bgcolor=gray!6]{bash}
dpkg -l                          # все пакеты и их состояние
dpkg -L <pkg>                    # файлы, установленные пакетом
dpkg -S <path>                   # какой пакет установил файл
dpkg -s <pkg>                    # подробная информация
dpkg -C                          # проверка состояния, поиск ошибок
\end{minted}

\section{Стандартные каталоги}

Пакеты в Debian/Ubuntu обычно ставят:
\begin{itemize}
  \item исполняемые файлы --- в \fpath{/usr/bin} и \fpath{/usr/sbin};
  \item библиотеки --- в \fpath{/usr/lib};
  \item документацию --- в \fpath{/usr/share/doc/<pkg>};
  \item конфигурацию --- в \fpath{/etc/<pkg>};
  \item метаданные --- в \fpath{/var/lib/dpkg/info}.
\end{itemize}

\cmd{dpkg -L} показывает реальный список для конкретного
пакета --- не полагайтесь на память.

\section{Типичные ошибки}
\begin{itemize}
  \item \emph{Unable to locate package} --- сначала
    \cmd{sudo apt update};
  \item Пакет установлен, но команда не находится ---
    посмотрите \cmd{dpkg -L}, какой бинарник поставил пакет;
  \item После удаления пакета остались файлы конфигурации ---
    это нормально; уберите \cmd{sudo apt purge <pkg>};
  \item Зависимости не разрешаются --- не используйте
    \cmd{dpkg -i} для пакетов с зависимостями;
  \item \cmd{apt} ругается на «held» пакеты --- посмотрите
    \cmd{apt-mark showhold}.
\end{itemize}

\section{Связь с остальным курсом}

Пакетная система нужна и для установки
\cmd{gcc}, \cmd{make}, \cmd{valgrind}
(\nameref{ch:13-c-compilation}, \nameref{ch:17-debugging}),
и для \cmd{git} (\nameref{ch:25-git}), и для \cmd{docker}
(\nameref{ch:19-docker-basics}). В Docker-образе пакетная
система --- тоже \cmd{apt} (если образ на Debian/Ubuntu),
см.\ \nameref{ch:20-dockerfile}.

\section{Задачи для самопроверки}
\begin{enumerate}
  \item Объясните разницу между \cmd{apt update} и
    \cmd{apt upgrade}. Что делает каждое?
  \item Чем \cmd{dpkg} отличается от \cmd{apt} в
    разрешении зависимостей?
  \item Как узнать, какие файлы принадлежат
    установленному пакету?
  \item Чем \cmd{remove} отличается от \cmd{purge}?
  \item Как установить локальный \fpath{.deb} так,
    чтобы зависимости разрешались автоматически?
\end{enumerate}

% Глава: 09-vim
% Статус: наполнено (первый прогон).
\chapter{Текстовый редактор Vim}
\label{ch:09-vim}

\section{Идея и моды}

\cmd{vim} (vi editor) --- клавиатурный редактор, установленный
практически в любой Linux-системе. Ключевая идея: \textbf{моды}.
Не всё время вы в режиме «печатаю текст»: есть отдельные режимы
для навигации, команд и ввода.

Vim работает в терминале вашей ВМ на \fpath{pglab.nsu.ru} --- это
основной способ править файлы на сервере без графического редактора.

\begin{center}
\begin{tabular}{lp{0.7\linewidth}}
\toprule
\textbf{Режим} & \textbf{Назначение} \\
\midrule
Normal   & навигация и вызов команд; режим по умолчанию \\
Insert   & ввод текста (аналог «обычного» редактора) \\
Visual   & выделение фрагмента для операций \\
Command  & строка команд в самом низу (\cmd{:}) для записи, выхода и т.п. \\
\bottomrule
\end{tabular}
\end{center}

Наиболее частая ошибка новичка --- попробовать печатать в режиме
Normal: текст не появляется. Чтобы начать печатать, нужно
войти в Insert (\cmd{i} --- вставить перед курсором, \cmd{a} ---
после, \cmd{o} --- новая строка снизу). Вернуться в Normal ---
нажать \cmd{Esc}.

\section{Базовая навигация (Normal mode)}

Перемещение без мыши:
\begin{minted}[bgcolor=gray!6]{bash}
h  j  k  l      # влево, вниз, вверх, вправо
w  b            # следующее / предыдущее слово
0  $            # начало / конец строки
gg  G           # начало / конец файла
Ctrl-d Ctrl-u   # полстраницы вниз / вверх
\end{minted}

Поиск и переход:
\begin{minted}[bgcolor=gray!6]{bash}
/pattern        # поиск вперёд от курсора
n  N            # следующее / предыдущее совпадение
\end{minted}

\section{Операции над текстом (Normal mode)}

Vim строит операции как \emph{глагол + объект}: \cmd{d} (delete),
\cmd{y} (yank/copy), \cmd{p} (paste), \cmd{c} (change) плюс
объект: \cmd{w} (word), \cmd{d} (line), \cmd{i}+пары (inside
parentheses/quotes), \cmd{t} (tag) и т.д.
\begin{minted}[bgcolor=gray!6]{bash}
dd         # удалить текущую строку
yy         # скопировать текущую строку
p          # вставить после курсора
u          # отменить
.          # повторить последнюю команду
dw         # удалить слово
ci"        # изменить внутри кавычек
\end{minted}

Операции можно умножать: \cmd{3dd} --- удалить 3 строки,
\cmd{5w} --- вперёд 5 слов.

\section{Запись и выход (Command mode)}

Внизу нажимаете \cmd{:} и вводите команду:
\begin{minted}[bgcolor=gray!6]{bash}
:w          # сохранить
:w new.txt  # сохранить как new.txt
:wq         # сохранить и выйти
:q          # выйти; не выйдет при несохранённых изменениях
:q!         # выйти, отбросив изменения
\end{minted}

Типовой сценарий: \cmd{Esc} для возврата в Normal, затем
\cmd{:wq} --- сохранить и выйти.

\section{Начальная настройка}

Vim читает \fpath{~/.vimrc} --- ваш конфиг.
Минимально полезный набор:
\begin{minted}[bgcolor=gray!6]{bash}
" show line numbers
set number
" 4 spaces, no tabs
set expandtab tabstop=4 shiftwidth=4
" highlight the current line
set cursorline
" sensible defaults for a modern system
set nocompatible
\end{minted}

Для работы с проектами, синтаксисом и автодополнением обычно
используют более развитые среды (Neovim, Emacs, IDE) --- но
навык базового Vim остаётся полезным на удалённых серверах,
в контейнерах и в rescue-режимах, где тяжёлый редактор
не поставить.

\section{Типичные ошибки}
\begin{itemize}
  \item Текст не печатается --- вы в Normal; нажмите \cmd{i}.
  \item \cmd{:q} не выходит --- есть несохранённые изменения;
    используйте \cmd{:q!} или сначала \cmd{:w}.
  \item Удалили случайно много строк --- \cmd{u} (undo) или
    \cmd{u} несколько раз; \cmd{U} --- восстановить строку целиком.
  \item Перепутали \cmd{d} (delete) и \cmd{y} (yank) ---
    данные потеряны? \cmd{p} вставит то, что было в буфере.
  \item Не нашли нужную строку --- \cmd{/pattern}, затем \cmd{n}.
\end{itemize}

\section{Связь с остальным курсом}

Vim удобен для правки конфигов
(\fpath{/etc/hosts}, \fpath{.bashrc}), Makefile
(\nameref{ch:14-make}), Dockerfile
(\nameref{ch:20-dockerfile}), C-кода
(\nameref{ch:13-c-compilation}). Внутри tmux
(\nameref{ch:10-tmux}) Vim --- основной редактор;
внутри контейнеров (\nameref{ch:19-docker-basics}) он часто
единственный доступный.

\section{Задачи для самопроверки}
\begin{enumerate}
  \item Объясните, чем отличаются \cmd{i}, \cmd{a}, \cmd{o}
    для входа в Insert.
  \item Как удалить текущую строку и как её восстановить?
  \item Как найти все вхождения слова и переключаться между
    ними?
  \item Чем \cmd{:q} отличается от \cmd{:q!}?
  \item Как в Vim перейти в конец файла, затем удалить
    последние три строки?
\end{enumerate}

% Глава: 10-tmux
% Статус: наполнено (первый прогон).
\chapter{Управление сессиями: tmux}
\label{ch:10-tmux}

\section{Зачем нужен tmux}

При работе по SSH (\nameref{ch:02-linux-access}) разрыв
соединения убивает все запущенные в терминале процессы.
\textbf{tmux} --- менеджер терминальных сессий: он позволяет
запустить программу в \emph{сессии}, отключиться, вернуться
позже (с другого компьютера) и продолжить работу --- процессы
останутся.

Для работы на сервере \fpath{pglab.nsu.ru} это особенно важно:
длинная компиляция или тесты не прервутся, даже если вы
случайно закрыли SSH-сессию.

Основные понятия:
\begin{center}
\begin{tabular}{lp{0.7\linewidth}}
\toprule
\textbf{Понятие} & \textbf{Что это} \\
\midrule
session & именованное пространство (несколько окон) \\
window  & «вкладка» внутри сессии \\
pane    & область внутри окна (деление экрана) \\
\bottomrule
\end{tabular}
\end{center}

\begin{figure}[htb]
  \centering
  \begin{tikzpicture}[node distance=1.2cm]
    \node[process, minimum width=6cm] (s) {session: \texttt{work}};
    \node[ext, below left=1.1cm and -0.5cm of s] (w1) {window: monitor};
    \node[ext, below right=1.1cm and -0.5cm of s] (w2) {window: editor};
    \node[data, below left=1.1cm and -0.3cm of w1] (p1) {pane: generator};
    \node[data, below right=1.1cm and -0.3cm of w1] (p2) {pane: tail -f};
    \node[data, below=1.1cm of w2] (p3) {pane: vim};
    \draw[flow] (s) -- (w1);
    \draw[flow] (s) -- (w2);
    \draw[flow] (w1) -- (p1);
    \draw[flow] (w1) -- (p2);
    \draw[flow] (w2) -- (p3);
  \end{tikzpicture}
  \caption{Иерархия tmux: сессия содержит окна, окно --- панели.}
  \label{fig:tmux}
\end{figure}

\section{Основное управление}

Большинство команд начинаются с \textbf{prefix} (по умолчанию
\cmd{Ctrl+b}), затем нажимается клавиша-команда. Ниже ---
без prefix там, где он не нужен.

\begin{minted}[bgcolor=gray!6]{bash}
tmux                      # новая сессия по умолчанию
tmux new -s work          # именованная сессия
tmux ls                   # список сессий
tmux attach -t work       # присоединиться к сессии
tmux kill-session -t work # удалить сессию
\end{minted}

Отключение от сессии (процессы продолжают работать):
нажмите \cmd{Ctrl+b}, отпустите, затем \cmd{d} (detach).

\section{Окна и панели (после prefix \cmd{Ctrl+b})}

\begin{center}
\begin{tabular}{lp{0.7\linewidth}}
\toprule
\textbf{Клавиша} & \textbf{Действие} \\
\midrule
\cmd{c}         & создать новое окно \\
\cmd{n} / \cmd{p} & следующее / предыдущее окно \\
\cmd{0}--\cmd{9} & перейти к окну по номеру \\
\cmd{\%}        & разделить окно вертикально \\
\cmd{\textbackslash} & разделить окно горизонтально \\
\cmd{h/j/k/l}  & переключение между панелями \\
\cmd{z}         & zoom: на весь экран / обратно \\
\cmd{x}         & закрыть текущую панель \\
\cmd{d}         & отключиться от сессии (detach) \\
\bottomrule
\end{tabular}
\end{center}

Типовая рабочая раскладка: в одной панели запущена долгая
программа, в другой --- наблюдение за логом (\cmd{tail -f}),
в третьем окне --- редактор.

\section{Настройка}

tmux читает \fpath{~/.tmux.conf}. Минимальный
полезный конфиг:
\begin{minted}[bgcolor=gray!6]{bash}
# использовать C-a как префикс вместо C-b
set -g prefix C-a
unbind C-b
bind C-a send-prefix
# нумеровать окна и панели с 1
set -g base-index 1
set -g pane-base-index 1
\end{minted}

\section{Типичные ошибки}
\begin{itemize}
  \item Нажали \cmd{Ctrl+b}, отпустили и ничего не произошло ---
    нужно нажать \emph{вторую} клавишу команды.
  \item Отключились и забыли имя сессии --- \cmd{tmux ls}.
  \item Сессия «умерла» --- её закрыли (kill) или сервер tmux
    завершился; detach от этого отличается.
  \item Не можете найти панель --- \cmd{Ctrl+b z} переключает
    zoom, \cmd{Ctrl+b q} показывает номера панелей.
\end{itemize}

\section{Связь с остальным курсом}

tmux --- стандарт для работы на сервере
(\nameref{ch:02-linux-access}) и в Docker
(\nameref{ch:19-docker-basics}): внутри контейнера можно
поднять tmux и иметь несколько панелей для наблюдения
за сервисами. Наблюдение за процессами в панелях ---
переход к \nameref{ch:12-processes}.

\section{Задачи для самопроверки}
\begin{enumerate}
  \item Чем detach от сессии отличается от kill сессии?
  \item Запустите сессию, создайте два окна и разделите одно
    окно на две панели. Как между ними переключаться?
  \item Как отключиться от сессии, не убивая процессы?
  \item Как переименовать или перейти к окну по номеру?
  \item Как вернуть панель из zoom?
\end{enumerate}

% Глава: 11-mc
% Статус: наполнено (первый прогон).
\chapter{Файловый менеджер: Midnight Commander}
\label{ch:11-mc}

\section{Зачем нужен}

\cmd{mc} (Midnight Commander) --- терминальный файловый
менеджер с «двухпанельным» интерфейсом. Полезен, когда
файловые операции проще сделать «глазами», чем командами:
сравнить две директории, переместить группу файлов,
просмотреть файл встроенным просмотрщиком.

Полезен, когда вы работаете с файлами на своей ВМ на
\fpath{pglab.nsu.ru} и нужен быстрый «визуальный» контроль.

\begin{figure}[htb]
  \centering
  \begin{tikzpicture}[node distance=1.4cm]
    \node[box, align=left, minimum width=4.5cm, minimum height=3.5cm] (left)
      {ЛЕВАЯ ПАНЕЛЬ\\
       \texttt{incoming/}\\
       \texttt{a.txt}\\
       \texttt{b.txt}};
    \node[box, align=left, right=2.5cm of left, minimum width=4.5cm,
          minimum height=3.5cm] (right)
      {ПРАВАЯ ПАНЕЛЬ\\
       \texttt{processed/}\\
       \texttt{a.txt}};
    \node[below=1.2cm of left, font=\small] {командная строка и меню внизу};
  \end{tikzpicture}
  \caption{Схема интерфейса Midnight Commander.}
  \label{fig:mc}
\end{figure}

\section{Навигация}

Внутри \cmd{mc}:
\begin{itemize}
  \item \cmd{Tab} --- переключить активную панель;
  \item стрелки --- перемещение по списку;
  \item \cmd{Enter} --- войти в каталог;
  \item \cmd{Backspace} / \cmd{Alt+u} --- вверх;
  \item \cmd{Ctrl+t} --- создать каталог;
  \item \cmd{Alt+s} --- показать свойства файла;
  \item \cmd{Alt+g} --- перейти по пути вручную;
  \item \cmd{Alt+z} / \cmd{Alt+x} --- сжать / распаковать архив;
  \item \cmd{F3} --- просмотр файла, \cmd{F4} --- открыть в редакторе.
\end{itemize}

\section{Файловые операции}

Выделяйте файлы клавишей \cmd{Insert} (или \cmd{Space}) и
используйте \cmd{F5} (копировать), \cmd{F6} (переместить/
переименовать), \cmd{F8} (удалить). Все операции
выполняются с подтверждением и показывают прогресс.

Полезные сочетания:
\begin{itemize}
  \item \cmd{F7} --- создать каталог;
  \item \cmd{F8} --- удалить;
  \item \cmd{F6} --- переместить или переименовать;
  \item \cmd{F5} --- копировать;
  \item \cmd{Ctrl+r} --- сравнить две директории (включая
    содержимое).
\end{itemize}

\section{Поиск и сортировка}

\cmd{Alt+s} --- свойства, \cmd{Ctrl+s} / \cmd{Alt+s} ---
поиск по имени; \cmd{Alt+s} --- переключить сортировку
(по имени, дате, размеру). \cmd{Ctrl+h} --- показать/
скрыть скрытые файлы (начинающиеся с \cmd{.}).

\section{Когда mc, а когда shell-команды}

\begin{center}
\begin{tabular}{lp{0.55\linewidth}}
\toprule
\textbf{Задача} & \textbf{Удобнее} \\
\midrule
одиночный файл & \cmd{cp} / \cmd{mv} \\
копировать дерево & \cmd{cp -r} \\
просмотр дерева & \cmd{mc} / \cmd{tree} \\
фильтр по имени & \cmd{find} + \cmd{grep} \\
сравнить две директории & \cmd{mc} (Ctrl+r) \\
пакетная обработка & shell-скрипт \\
\bottomrule
\end{tabular}
\end{center}

\cmd{mc} не заменяет shell-команды --- это другой интерфейс
к тем же операциям. Для повторяющихся задач скрипт на
оболочке (\nameref{ch:27-scripting-ide}) масштабируете,
а \cmd{mc} --- для разовых манипуляций.

\section{Типичные ошибки}
\begin{itemize}
  \item Не видите скрытые файлы --- включите
    \cmd{Ctrl+h} (или настройку «Show hidden files»).
  \item Переместили файл не туда --- посмотрите, какая
    панель была \emph{активной} (обведённая рамкой);
    операция выполняется из активной в неактивную.
  \item Не выходите из \cmd{mc}, нажав \cmd{q} ---
    используйте \cmd{F10} или \cmd{Ctrl+q}.
\end{itemize}

\section{Связь с остальным курсом}

\cmd{mc} удобен для навигации по структуре, созданной
пакетами (\nameref{ch:08-packages}) или Docker
(\nameref{ch:21-docker-data}): в контейнере, где нет
\cmd{tree}, \cmd{mc} часто единственный графический
по терминалу интерфейс для файлов.

\section{Задачи для самопроверки}
\begin{enumerate}
  \item Как переключаться между левой и правой панелью?
  \item Как создать в активной панели каталог
    и перейти в него?
  \item Как скопировать несколько выделенных файлов
    в другую директорию?
  \item Как сравнить две директории на различия?
  \item Чем работа в \cmd{mc} отличается от \cmd{find}
    + \cmd{grep} по гибкости?
\end{enumerate}

% Глава: 12-processes
% Статус: наполнено (первый прогон).
\chapter{Процессы и ресурсы}
\label{ch:12-processes}

\section{Что такое процесс}

\textbf{Процесс} --- это программа в состоянии исполнения:
свой набор инструкций, память, открытые файлы, потоки.
Каждый процесс получает уникальный идентификатор \textbf{PID}
(Process ID) и работает от имени пользователя
(\nameref{ch:05-permissions}).

В этой главе процессы --- это то, что запущено на вашей ВМ:
компиляторы, тесты, контейнеры на сервере \fpath{pglab.nsu.ru}.

Процесс имеет \textbf{состояния} (видны в \cmd{ps}):
\begin{itemize}
  \item \cmd{R} --- running / runnable (выполняется или ждёт CPU);
  \item \cmd{S} --- sleeping (ждёт событие: ввод, диск, сеть);
  \item \cmd{Z} --- zombie (завершился, но родитель ещё не
    «забрал» результат).
\end{itemize}

\section{Иерархия процессов}

Все процессы, кроме первого (\cmd{init} / \cmd{systemd}),
являются чьими-то дочерними. Родительский идентификатор ---
\textbf{PPID}. Когда оболочка запускает команду
(\nameref{ch:03-shell}), команда --- дочерний процесс
оболочки.

\begin{figure}[htb]
  \centering
  \begin{tikzpicture}[node distance=1.3cm]
    \node[process] (init) {systemd (PID 1)};
    \node[process, below left=1.2cm and -0.4cm of init] (sshd) {sshd};
    \node[process, below right=1.2cm and -0.4cm of init] (cron) {cron};
    \node[process, below left=1.2cm and -0.4cm of sshd] (bash) {bash};
    \node[process, below right=1.2cm and -0.4cm of sshd] (vim) {vim};
    \draw[flow] (init) -- (sshd);
    \draw[flow] (init) -- (cron);
    \draw[flow] (sshd) -- (bash);
    \draw[flow] (sshd) -- (vim);
  \end{tikzpicture}
  \caption{Дерево процессов: каждый процесс имеет родителя.}
  \label{fig:proc-tree}
\end{figure}

\section{Просмотр процессов}

\begin{minted}[bgcolor=gray!6]{bash}
ps aux                 # все процессы: пользовательские и системные
ps -ef                 # полный формат, с PPID
ps -u <user>           # процессы одного пользователя
ps aux | grep <name>   # найти процесс по именем
top                    # «живой» просмотр с обновлением
htop                   # интерактивно, с цветом, удобнее читать
\end{minted}

Колонки \cmd{ps aux}, которые стоит читать: \textbf{USER},
\textbf{PID}, \textbf{\%CPU}, \textbf{\%MEM},
\textbf{COMMAND}. В \cmd{top}/\cmd{htop} сверху видна
сводка: загрузка CPU, свободная память, количество
процессов.

\section{Сигналы и завершение процессов}

Процессу можно отправить \textbf{сигнал} --- системное
уведомление. Важнейшие:

\begin{center}
\begin{tabular}{clp{0.5\linewidth}}
\toprule
\textbf{Сигнал} & \textbf{Номер} & \textbf{Действие} \\
\midrule
SIGHUP  & 1  & «отключён терминал»; процесс обычно завершается \\
SIGINT  & 2  & прерывание (обычно \cmd{Ctrl+C} в терминале) \\
SIGQUIT & 3  & завершение с дампом (редко) \\
SIGTERM & 15 & вежливая просьба завершиться (по умолчанию) \\
SIGKILL & 9  & принудительное завершение (процесс не может перехватить) \\
\bottomrule
\end{tabular}
\end{center}

\begin{minted}[bgcolor=gray!6]{bash}
kill 1234              # SIGTERM для PID 1234
kill -9 1234           # SIGKILL, принудительно
kill -TERM <pid>       # явно
pkill <name>           # по имени процесса
pkill -f "python app"  # по полной командной строке
\end{minted}

Правило: сначала \cmd{SIGTERM}, и только если процесс не
реагирует --- \cmd{SIGKILL}. Вежливое завершение даёт
процессу закрыть файлы и освободить ресурсы.

\section{Фоновые и передние задачи}

Запустить процесс в фоне --- добавить \cmd{\&} в конец
команды:
\begin{minted}[bgcolor=gray!6]{bash}
long_job.sh &        # запустить в фоне
jobs               # фоновые задания этой оболочки
fg                 # вернуть последнее задание на передний план
fg %2              # вернуть задание №2 на передний план
bg                 # возобновить остановленное задание в фоне
Ctrl+Z             # остановить текущее задание
\end{minted}

\section{Ресурсы: CPU и память}

Linux выделяет процессу ресурсы по запросу: долю CPU
(время процессора) и виртуальную память
(\nameref{ch:16-virtual-memory}). Мониторинг:
\begin{minted}[bgcolor=gray!6]{bash}
top                  # загрузка по CPU и по процессам
htop                 # интерактивно
vmstat 1             # системная статистика раз в секунду
free -h              # сводка по памяти
\end{minted}

\section{Типичные ошибки}
\begin{itemize}
  \item Процесс «завис» --- посмотрите в \cmd{top}, не
    использует ли он 100\% CPU; иногда достаточно
    \cmd{kill} (SIGTERM).
  \item \cmd{kill 9} не работает --- правильно
    \cmd{kill -9}; без знака минус это PID 9, а не флаг.
  \item Задача в фоне «потерялась» --- \cmd{jobs},
    затем \cmd{fg} или \cmd{bg}.
  \item После \cmd{Ctrl+C} сервер не запустить снова ---
    убедитесь, что старый процесс реально завершён
    (\cmd{ps aux | grep ...}).
\end{itemize}

\section{Связь с остальным курсом}

Процессы --- центральное понятие ОС. Потоки вывода
(\nameref{ch:07-text-pipelines}) принадлежат процессам;
запуск и остановка контейнеров
(\nameref{ch:19-docker-basics}) --- это запуск и
остановка процессов в изолированном окружении.
Взаимодействие процессов между собой ---
\nameref{ch:18-ipc}; распределение памяти ---
\nameref{ch:16-virtual-memory}.

\section{Задачи для самопроверки}
\begin{enumerate}
  \item Как узнать PID запущенного процесса?
  \item Чем \cmd{SIGTERM} отличается от \cmd{SIGKILL}?
    Почему второй «невозможно перехватить»?
  \item Запустите долгую команду в фоне, остановите её
    (\cmd{Ctrl+Z}), затем возобновите в фоне.
  \item Как увидеть, какие процессы принадлежат текущему
    пользователю?
  \item Что такое zombie-процесс и почему он возникает?
\end{enumerate}


\part{Разработка на C и сборка}
% Глава: 13-c-compilation
% Статус: наполнено (первый прогон).
\chapter{Язык C и компиляция}
\label{ch:13-c-compilation}

\section{Язык C: зачем и что важно}

C --- низкий, близкий к машинному коду язык: предсказуемые
свойства, отсутствие скрытых механизмов, прямой доступ к
памяти. Именно поэтому ядро Linux и большинство системных
инструментов написаны на C. Для курса важно понимать:
как работает компилятор, как линковщик собирает программу,
как устроена память.

Компилируем на сервере: \cmd{gcc} запускается на вашей ВМ
\fpath{pglab.nsu.ru}, и все артефакты сборки тоже остаются на ней.

\section{Компилятор и интерпретатор}

\textbf{Компилятор} (C) переводит весь исходный код в
машинный код \emph{до} запуска; результат --- исполняемый
файл. \textbf{Интерпретатор} (Python, Ruby, большинство
JS) выполняет исходный код «на лету». Компиляция даёт
быстроту выполнения, интерпретация --- простоту запуска.

\section{Этапы компиляции C}

Командой \cmd{gcc} вы обычно делаете всё сразу, но внутри
несколько этапов:

\begin{figure}[htb]
  \centering
  \begin{tikzpicture}[node distance=1.3cm]
    \node[data] (src) {source.c};
    \node[process, right=of src] (pp) {preprocess\\[0.2em] (gcc -E)};
    \node[process, right=of pp] (asm) {assemble\\[0.2em] to .o};
    \node[process, right=of asm] (lnk) {link\\[0.2em] (ld)};
    \node[data, right=of lnk] (bin) {binary};
    \draw[flow] (src) -- (pp);
    \draw[flow] (pp) -- (asm);
    \draw[flow] (asm) -- (lnk);
    \draw[flow] (lnk) -- (bin);
  \end{tikzpicture}
  \caption{Этапы: препроцессор, компиляция в объект, линковка.}
  \label{fig:c-pipeline}
\end{figure}

\begin{itemize}
  \item \textbf{Preprocessing} (\cmd{gcc -E}) --- раскрытие
    \cmd{\#include}, \cmd{\#define}, макросов;
  \item \textbf{Compilation} --- перевод в assembly;
  \item \textbf{Assembly} --- перевод assembly в машинный
    код объектного файла \fpath{.o};
  \item \textbf{Linking} --- сборка объектных файлов и
    библиотек в один исполняемый файл.
\end{itemize}

\section{Типовой сценарий}

\begin{minted}[bgcolor=gray!6]{c}
// hello.c
#include <stdio.h>

int main(void) {
    int n = 42;
    printf("answer: %d\n", n);
    return 0;
}
\end{minted}

\begin{minted}[bgcolor=gray!6]{bash}
gcc -Wall -Wextra -o hello hello.c   # за один раз: исходник -> бинарник
./hello                              # запустить
gcc -c hello.c                       # только компиляция: hello.c -> hello.o
ld hello.o -o hello -lc              # собрать вручную, нужна библиотека
\end{minted}

\section{Флаги \cmd{gcc}}

\begin{center}
\begin{tabular}{lp{0.65\linewidth}}
\toprule
\textbf{Флаг} & \textbf{Назначение} \\
\midrule
\cmd{-c}        & только компиляция в объектный файл, без линковки \\
\cmd{-o name}   & имя выходного файла \\
\cmd{-I dir}    & каталог для поиска заголовков (\cmd{\#include}) \\
\cmd{-L dir}    & каталог для поиска библиотек на этапе линковки \\
\cmd{-l lib}    & подключить библиотеку (например \cmd{-lm} --- math) \\
\cmd{-Wall} \cmd{-Wextra} & включение большинства предупреждений \\
\cmd{-g}        & отладочная информация (для gdb, valgrind) \\
\cmd{-O0 -O1 -O2 -O3} & уровни оптимизации \\
\cmd{-std=c11}  & выбрать стандарт языка \\
\bottomrule
\end{tabular}
\end{center}

\begin{tipbox}
Всегда компилируйте с \cmd{-Wall -Wextra}: предупреждения
компилятора ловят типовые ошибки (неиспользуемые переменные,
потерянные значения) раньше, чем рантайм.
\end{tipbox}

\section{Заголовочные файлы и include-пути}

C делит код на \textbf{модули} (\fpath{.c} файлы) с
публичным интерфейсом в \fpath{.h} заголовках. \cmd{\#include
"file.h"} ищет в текущем каталоге, \cmd{\#include <file.h>} ---
в системных каталогах. Если заголовки лежат не в
системных каталогах, укажите их через \cmd{-I}.

Типовой проект:
\begin{minted}[bgcolor=gray!6]{bash}
project/
  src/main.c
  src/util.c
  src/util.h
  Makefile
\end{minted}

\begin{minted}[bgcolor=gray!6]{bash}
gcc -Wall -Wextra -Isrc -o app src/main.c src/util.c
\end{minted}

\section{Библиотеки}

\textbf{Статическая} (\fpath{libfoo.a}) --- код библиотеки
встраивается в ваш бинарник; бинарник самодостаточен.
\textbf{Динамическая} (\fpath{libfoo.so}) --- общий код
на диске; бинарник ссылается на него во время выполнения.
Динамические меньше по размеру и позволяют обновлять
библиотеку без пересборки.

\begin{minted}[bgcolor=gray!6]{bash}
gcc -o prog main.c -L./lib -lfoo      # линковка с libfoo
ldd prog                              # показать используемые общие библиотеки
\end{minted}

\section{Типичные ошибки}
\begin{itemize}
  \item \emph{undefined reference to ...} --- линковщик
    не нашёл функцию; проверьте, все ли \fpath{.c}
    переданы и не забыта ли библиотека.
  \item \emph{implicit declaration of function} --- забыли
    \cmd{\#include} или опечатались в имени; при
    \cmd{-Wall} это предупреждение, но лучше править.
  \item Выходной файл не запускается --- не тот процессор
    или отсутствует библиотека; \cmd{ldd} покажет, что
    не хватает.
  \item Код «работает, но странно» --- чаще всего это
    обращение к неинициализированной памяти; см.
    \nameref{ch:17-debugging}.
\end{itemize}

\section{Связь с остальным курсом}

Сборка проекта через \cmd{make} --- \nameref{ch:14-make};
тестирование и оптимизация --- \nameref{ch:15-testing-optimization};
память и как её распределяет ОС ---
\nameref{ch:16-virtual-memory}; отладка и поиск утечек ---
\nameref{ch:17-debugging}. Docker-образы с C-проектами
часто включают \cmd{gcc} и \cmd{make}
(\nameref{ch:20-dockerfile}).

\section{Задачи для самопроверки}
\begin{enumerate}
  \item Объясните четыре этапа компиляции C. Что
    генерирует каждый из них?
  \item Чем \cmd{gcc -c} отличается от \cmd{gcc}?
  \item Как скомпилировать проект из нескольких
    \fpath{.c} с заголовками в подкаталоге?
  \item В чём разница между статической и динамической
    библиотекой?
  \item Какие флаги вы бы использовали для отладочной
    и для оптимизированной сборки?
\end{enumerate}

% Глава: 14-make
% Статус: наполнено (первый прогон).
\chapter{Система сборки: make}
\label{ch:14-make}

\section{Зачем нужен make}

Когда проект вырастает за один \cmd{gcc}, ручная сборка
становится ненадёжной: легко забыть файл или пересобрать
то, что не менялось. \textbf{make} --- система сборки:
он смотрит на \emph{зависимости} (что из чего) и
пересобирает \textbf{только изменившееся}.

Выполняем сборку на своей ВМ: утилита \cmd{make} читает
\fpath{Makefile} на сервере \fpath{pglab.nsu.ru} и пересобирает нужное.

\section{Анатомия Makefile}

Файл \fpath{Makefile} состоит из \textbf{правил}
(rules):
\begin{minted}[bgcolor=gray!6]{Makefile}
target: prerequisites
    recipe command
\end{minted}

\begin{itemize}
  \item \textbf{target} --- что строим (файл или имя);
  \item \textbf{prerequisites} --- от чего зависит;
  \item \textbf{recipe} --- команды (всегда с \emph{табуляцией},
    не с пробелами!).
\end{itemize}

\begin{warnbox}
В листингах этой книги отступ рецепта показан \textbf{четырьмя
пробелами} --- табуляцию внутри блока напечатать нельзя. В настоящем
\file{Makefile} перед командой рецепта \textbf{обязательно стоит
одна таб}: make с пробелами в начале строки не примет и выдаст
ошибку \emph{missing separator}.
\end{warnbox}

Минимальный Makefile для \fpath{hello.c}:
\begin{minted}[bgcolor=gray!6]{Makefile}
CC = gcc
CFLAGS = -Wall -Wextra -g

hello: hello.c
    $(CC) $(CFLAGS) -o hello hello.c

clean:
    rm -f hello
\end{minted}

\begin{warnbox}
Рецепты в Makefile \textbf{обязаны} начинаться с табуляции.
Пробелы --- синтаксическая ошибка. При копировании из
редактора проверяйте, что это именно табы.
\end{warnbox}

\section{Переменные и автоматические переменные}

Обычные переменные задаются как \cmd{CC = gcc}.
Автоматические переменные зависят от текущего правила:
\begin{center}
\begin{tabular}{lp{0.6\linewidth}}
\toprule
\textbf{Переменная} & \textbf{Значение} \\
\midrule
\cmd{$@}   & целевой файл текущего правила \\
\cmd{$<}   & первый из зависимостей \\
\cmd{$^}   & все зависимости \\
\cmd{$(CC)} & значение переменной \\
\bottomrule
\end{tabular}
\end{center}

Использование \cmd{\$@} и \cmd{\$<} укорачивает правила:
\begin{minted}[bgcolor=gray!6]{Makefile}
app: main.o util.o
    $(CC) -o $@ $^
\end{minted}

\section{Паттерн-правила}

Паттерн-правило описывает, как строить \emph{любой} файл
по шаблону. Классический пример: из \fpath{.c} в \fpath{.o}:
\begin{minted}[bgcolor=gray!6]{Makefile}
%.o: %.c
    $(CC) $(CFLAGS) -c $< -o $@
\end{minted}

\cmd{\%} --- заглушка, подставляющая имя.

\section{Типовой Makefile для проекта}

\begin{minted}[bgcolor=gray!6]{Makefile}
CC      = gcc
CFLAGS  = -Wall -Wextra -g
SRCDIR  = src
OBJDIR  = obj
SRC     = $(wildcard $(SRCDIR)/*.c)
OBJ     = $(SRC:$(SRCDIR)/%.c=$(OBJDIR)/%.o)
DEPS    = $(OBJ:.o=.d)

all: app

app: $(OBJ)
    $(CC) $(CFLAGS) -o $@ $^

$(OBJDIR)/%.o: $(SRCDIR)/%.c | $(OBJDIR)
    $(CC) $(CFLAGS) -MMD -MP -c $< -o $@

$(OBJDIR):
    mkdir -p $(OBJDIR)

-include $(DEPS)

clean:
    rm -rf $(OBJDIR) app

.PHONY: all clean
\end{minted}

Обратите внимание: \cmd{.PHONY} говорит make, что
\cmd{all} и \cmd{clean} --- не файлы; это важно, если
на диске появятся файлы с такими именами.

\section{Полезные флаги \cmd{make}}

\begin{minted}[bgcolor=gray!6]{bash}
make                 # целевая по умолчанию, обычно all
make clean           # выполнить целевую clean
make -n              # вывести команды без выполнения (dry run)
make -j4             # запускать до 4 задач параллельно
make -f other.mk     # использовать другой makefile
make CC=clang        # переопределить переменную в командной строке
\end{minted}

\section{Типичные ошибки}
\begin{itemize}
  \item \emph{missing separator} --- в рецепте пробелы
    вместо табуляции;
  \item Циклическая зависимость --- make останавливается
    с ошибкой;
  \item Пересборка каждый раз --- пропущена зависимость
    (например, заголовочный файл); используйте
    \cmd{-MMD -MP} для автоматического отслеживания;
  \item \cmd{make clean} делает что-то лишнее --- проверьте,
    что в рецепте только нужные команды.
\end{itemize}

\section{Связь с остальным курсом}

\cmd{make} --- стандарт сборки C-проектов
(\nameref{ch:13-c-compilation}) и в Docker-образах
(\nameref{ch:20-dockerfile}). Понимание зависимостей
пригодится и в скриптах CI
(\nameref{ch:26-ci}).

\section{Задачи для самопроверки}
\begin{enumerate}
  \item Из чего состоит правило в Makefile?
  \item Зачем \cmd{.PHONY}?
  \item Чем \cmd{\$@} отличается от \cmd{\$<}?
  \item Что делает \cmd{make -n}?
  \item Как заставить make пересобирать проект при
    изменении заголовочного файла?
\end{enumerate}

% Глава: 15-testing-optimization
% Статус: наполнено (первый прогон).
\chapter{Тестирование и оптимизация}
\label{ch:15-testing-optimization}

\section{Зачем тестировать}

Код, который «работает сейчас», не значит, что будет
работать после правок. \textbf{Тест} --- это воспроизводимый
проверочный сценарий: «при входе $X$ программа должна
дать $Y$». Набор тестов --- живая документация и
страховка при рефакторинге.

Запускаем тесты на своей ВМ на \fpath{pglab.nsu.ru}: так они
выполняются в той же среде, где всё собиралось.

\section{Что такое тест-кейс}

Тест-кейс определяет: \textbf{входные данные},
\textbf{ожидаемый результат} и способ
сравнения. Хороший тест:
\begin{itemize}
  \item проверяет \emph{одну} функцию/сценарий;
  \item не зависит от порядка и состояния других тестов;
  \item выполняется быстро (миллисекунды);
  \item сообщает, \emph{что именно} не так при провале.
\end{itemize}

\section{\cmd{assert} в C}

Простейший механизм --- \cmd{assert}:
\begin{minted}[bgcolor=gray!6]{c}
#include <assert.h>

int add(int a, int b) { return a + b; }

int main(void) {
    assert(add(2, 3) == 5);
    assert(add(10, 20) == 30);
    printf("all assertions passed\n");
    return 0;
}
\end{minted}

\cmd{assert(условие)} при ложном условии выводит сообщение
и завершает программу (через \cmd{abort}). В
оптимизированной сборке с \cmd{-DNDEBUG} все
\cmd{assert} отключаются.

\begin{tipbox}
\cmd{assert} --- не замена полному тест-фреймворку: он
не умеет «продолжить после провала» и собирать статистику.
Для серьёзных проектов используют фреймворки (Check,
CUTL, Unity, CTest), но идея --- та же.
\end{tipbox}

\section{Организация тестов}

Отдельный файл \fpath{tests/test_main.c} подключает
функции из основного кода и вызывает тесты:
\begin{minted}[bgcolor=gray!6]{bash}
gcc -Wall -Wextra -g -o test_app src/util.c tests/test_main.c
./test_app
\end{minted}

В Makefile удобно завести отдельный тег:
\begin{minted}[bgcolor=gray!6]{Makefile}
test: test_app
    ./test_app

test_app: $(OBJ) tests/test_main.o
    $(CC) -o $@ $^
\end{minted}

\section{Замер времени выполнения}

\begin{minted}[bgcolor=gray!6]{bash}
time ./app            # встроенная команда: user/sys/real
/usr/bin/time -v ./app   # подробный отчёт
\end{minted}

Ключевые метрики:
\begin{itemize}
  \item \textbf{real} --- «по часам» (wall clock);
  \item \textbf{user} --- время процессора в пользовательском
    режиме;
  \item \textbf{sys} --- время в режиме ядра.
\end{itemize}

Для сравнения версий алгоритмов измеряйте \emph{real} на
однаковых данных.

\section{Оптимизация}

Порядок: сначала убедитесь, что код \emph{правильный},
затем \emph{быстрый}. Уровни оптимизации \cmd{gcc}:
\begin{center}
\begin{tabular}{lp{0.6\linewidth}}
\toprule
\textbf{Флаг} & \textbf{Что делает} \\
\midrule
\cmd{-O0} & отладочная сборка без оптимизаций (по умолчанию) \\
\cmd{-O1} & лёгкие, безопасные оптимизации \\
\cmd{-O2} & основные оптимизации (хороший баланс) \\
\cmd{-O3} & агрессивные: развёртка, векторизация \\
\bottomrule
\end{tabular}
\end{center}

Оптимизация \textbf{ломает отладку}: номера строк и
значения переменных могут «гулять». Для отладки ---
\cmd{-O0 -g}, для релиза --- \cmd{-O2}.

\section{Типичные ошибки}
\begin{itemize}
  \item Тесты не проверяют граничные случаи (0, отрицательные,
    максимальные значения).
  \item \cmd{assert} в «боевом» коде, а не в тестах ---
    при \cmd{-DNDEBUG} проверки пропадают.
  \item Оптимизируют прежде, чем измерить --- без
    профиля оптимизация часто бессмысленна.
  \item Сравнивают время без \cmd{-O2} и на разных
    данных --- результат неинформативен.
\end{itemize}

\section{Связь с остальным курсом}

Тесты запускаются через \cmd{make}
(\nameref{ch:14-make}), оптимизация и отладочная
информация --- флаги \cmd{gcc}
(\nameref{ch:13-c-compilation}), а поиск утечек
и некорректного доступа к памяти ---
\nameref{ch:17-debugging}. Автоматический запуск
тестов --- основа CI (\nameref{ch:26-ci}).

\section{Задачи для самопроверки}
\begin{enumerate}
  \item Что такое тест-кейс? Какие у него компоненты?
  \item Чем \cmd{assert} отличается от полноценного
    тест-фреймворка?
  \item Как измерить время выполнения программы и
    какие метрики важны?
  \item Почему для отладки используют \cmd{-O0 -g},
    а для релиза --- \cmd{-O2}?
  \item Приведите пример граничного случая, который
    стоит проверить в тесте для функции сортировки.
\end{enumerate}

% Глава: 16-virtual-memory
% Статус: наполнено (первый прогон).
\chapter{Виртуальная память}
\label{ch:16-virtual-memory}

\section{Зачем нужна виртуальная память}

Каждому процессу хотелось бы видеть «свою» память, начиная
с адреса 0, и не сталкиваться с другими процессами.
\textbf{Виртуальная память} --- это механизм ОС, который
даёт каждому процессу \emph{отдельное адресное пространство},
а ядро переводит виртуальные адреса в физические.

Это же самое происходит на вашей ВМ: каждый процесс,
запущенный на \fpath{pglab.nsu.ru}, живёт в собственной
виртуальной памяти.

Выгоды:
\begin{itemize}
  \item \textbf{Изоляция} --- сбой в одном процессе не
    «ломает» соседний;
  \item \textbf{Больше памяти, чем в RAM} --- неиспользуемые
    страницы выгружаются на диск (\emph{swap});
  \item \textbf{Общий код} --- один и тот же текст библиотеки
    может жить в RAM один раз и «видеться» многим процессам.
\end{itemize}

\section{Строение адресного пространства процесса}

У процесса есть несколько областей:

\begin{figure}[htb]
  \centering
  \begin{tikzpicture}[node distance=0.6cm]
    \node[box, minimum width=7cm, minimum height=1cm] (stack) {stack (растёт вниз)};
    \node[box, minimum width=7cm, minimum height=1cm, below=of stack] (heap) {heap (malloc, растёт вверх)};
    \node[box, minimum width=7cm, minimum height=1.6cm, below=of heap] (data) {data / bss\\ (глобальные переменные)};
    \node[box, minimum width=7cm, minimum height=1cm, below=of data] (text) {text (код программы)};
  \end{tikzpicture}
  \caption{Основные области адресного пространства процесса.}
  \label{fig:mem-layout}
\end{figure}

\begin{itemize}
  \item \textbf{stack} --- локальные переменные и адреса
    возврата; управляется автоматически, ограничен
    (обычно несколько мегабайт);
  \item \textbf{heap} --- память, запрашиваемая через
    \cmd{malloc}/\cmd{new}; управляется программой;
  \item \textbf{data/bss} --- глобальные и статические
    переменные;
  \item \textbf{text} --- код программы (только чтение).
\end{itemize}

\section{Страницы и таблица страниц}

Память разбита на \textbf{страницы} (обычно 4 КБ).
Ядро хранит для процесса \textbf{таблицу страниц}:
соответствие «виртуальная страница $\to$ физическая
страница (или файл на диске)». При обращении к адресу
аппарат (\emph{MMU}/TLB) через таблицу переводит адрес;
если страницы нет в RAM --- происходит \textbf{page fault},
ядро подгружает её.

\begin{figure}[htb]
  \centering
  \begin{tikzpicture}[node distance=1.4cm]
    \node[process] (cpu) {CPU\\[0.2em] виртуальный адрес};
    \node[process, below=of cpu] (mmu) {MMU / TLB\\[0.2em] перевод адреса};
    \node[process, below=of mmu] (pt) {таблица страниц};
    \node[storage, below left=1.3cm and -0.6cm of pt] (ram) {RAM};
    \node[storage, below right=1.3cm and -0.6cm of pt] (disk) {диск (swap)};
    \draw[flow, <->] (cpu) -- (mmu);
    \draw[flow, <->] (mmu) -- (pt);
    \draw[flow, <->] (pt) -- (ram);
    \draw[flow, <->] (pt) -- (disk);
  \end{tikzpicture}
  \caption{Перевод виртуальных адресов в физические.}
  \label{fig:vm}
\end{figure}

\section{\cmd{mmap} и общая память}

Функция \cmd{mmap} позволяет «пришить» файл или
анонимную область к адресному пространству процесса.
Через неё реализованы и загрузка библиотек, и
\textbf{общая память} между процессами (см.\
\nameref{ch:18-ipc}): два процесса могут видеть одну
и ту же физическую страницу.

\section{Связь с процессами и Docker}

Каждый процесс (\nameref{ch:12-processes}) имеет своё
адресное пространство; когда процесс вызывается из
оболочки, у него отдельный стек и куча. Контейнеры
Docker (\nameref{ch:19-docker-basics}) разделяют
\emph{имена} и ресурсы, но виртуальная память процессов
внутри контейнера работает по тем же правилам;
ограничения по RAM на контейнер --- через cgroups
(\nameref{ch:23-docker-resources}).

\section{Типичные проблемы, связанные с памятью}

\begin{itemize}
  \item \textbf{Утечка} --- \cmd{malloc} без \cmd{free};
    память накапливается, процесс «раздувается»;
  \item \textbf{Переполнение стека} --- глубокая
    рекурсия или огромный локальный массив;
  \item \textbf{Use-after-free} --- обращение к памяти
    после \cmd{free};
  \item \textbf{Дangling pointer} --- указатель на
    объект, который уже уничтожен.
\end{itemize}

Всё это ищут инструментами из \nameref{ch:17-debugging}.

\section{Задачи для самопроверки}
\begin{enumerate}
  \item Зачем нужна виртуальная память? Назовите
    три выгоды.
  \item Перечислите основные области адресного
    пространства и что в них хранится.
  \item Что такое страница и таблица страниц?
  \item Что происходит при page fault?
  \item Чем heap отличается от stack по управлению
    и размеру?
\end{enumerate}

% Глава: 17-debugging
% Статус: наполнено (первый прогон).
\chapter{Отладка: gdb и Valgrind}
\label{ch:17-debugging}

\section{Два уровня отладки}

Когда программа «ведёт себя странно», есть два уровня:
\begin{itemize}
  \item \textbf{Логическая ошибка} --- код делает не то,
    что задумано; ищут чтением, принтами, stepping
    в отладчике;
  \item \textbf{Некорректный доступ к памяти} --- код
    обращается к чужой/свободной памяти; ищут
    специализированными инструментами (Valgrind).
\end{itemize}

Отладку ведём на своей ВМ: утилиты \cmd{gdb} и
\cmd{valgrind} работают на сервере \fpath{pglab.nsu.ru}, видя
процесс там, где он запущен.

Для обоих важно собирать с \cmd{-g} (отладочная
информация) и без агрессивной оптимизации
(\cmd{-O0}), чтобы номера строк и значения совпадали
с исходником.

\section{gdb: пошаговая отладка}

\cmd{gdb} --- стандартный отладчик GNU. Типовой цикл:
\begin{minted}[bgcolor=gray!6]{bash}
gcc -Wall -Wextra -g -o app app.c
gdb ./app
\end{minted}

Основные команды (внутри gdb):
\begin{center}
\begin{tabular}{lp{0.6\linewidth}}
\toprule
\textbf{Команда} & \textbf{Действие} \\
\midrule
\cmd{run}      & запустить программу \\
\cmd{break main} & поставить точку останова в main \\
\cmd{break file.c:42} & точка останова по строке \\
\cmd{next}     & следующий шаг (не входя в функции) \\
\cmd{step}     & шаг с заходом в функции \\
\cmd{continue}& продолжить до следующей точки \\
\cmd{print x}  & вывести значение переменной \\
\cmd{backtrace}& показать стек вызовов \\
\cmd{quit}     & выйти из gdb \\
\bottomrule
\end{tabular}
\end{center}

Полезно: \cmd{info locals} --- локальные переменные,
\cmd{watch var} --- остановка при изменении переменной.

\section{Valgrind: поиск проблем с памятью}

\cmd{valgrind} --- инструмент статического анализа
памяти во время выполнения. \cmd{memcheck} ---
наиболее полезный режим:
\begin{minted}[bgcolor=gray!6]{bash}
valgrind --leak-check=full ./app
\end{minted}

Отчёт Valgrind показывает:
\begin{itemize}
  \item \textbf{definitely lost} --- точно утекла память
    (нет ни одного указателя на неё);
  \item \textbf{indirectly lost} --- недоступна через
    другую утечку;
  \item \textbf{still reachable} --- блок жив, но нет
    указателей (обычно нормально на выходе);
  \item \textbf{invalid read/write} --- обращение за
    границами выделенного блока.
\end{itemize}

\begin{tipbox}
Запускайте Valgrind на \emph{отладочной} сборке
(\cmd{-g -O0}): он сам отключает оптимизации, но
с вашими флагами отчёт будет точнее.
\end{tipbox}

\section{Типовые ошибки памяти}

\begin{center}
\begin{tabular}{lp{0.6\linewidth}}
\toprule
\textbf{Ошибка} & \textbf{Что это} \\
\midrule
Buffer overflow   & запись за конец массива/строки \\
Use-after-free    & обращение к памяти после free \\
Dangling pointer  & указатель на уничтоженный объект \\
Double free       & дважды вызвать free на один блок \\
Off-by-one        & индекс на единицу не так (0..n вместо 0..n-1) \\
\bottomrule
\end{tabular}
\end{center}

Все эти ошибки в C «не падают немедленно»: программа
может работать нормально, а потом дать странный
результат или упасть. Valgrind находит их
детерминированно.

\section{Процесс отладки}

\begin{figure}[htb]
  \centering
  \begin{tikzpicture}[node distance=1.4cm]
    \node[data] (bug) {баг};
    \node[process, below=of bug] (repro) {воспроизвести};
    \node[process, below=of repro] (isolate) {изоляция (бинарный поиск, принты)};
    \node[process, below=of isolate] (gdb) {gdb / valgrind};
    \node[process, below=of gdb] (fix) {фикс + тест};
    \draw[flow] (bug) -- (repro);
    \draw[flow] (repro) -- (isolate);
    \draw[flow] (isolate) -- (gdb);
    \draw[flow] (gdb) -- (fix);
  \end{tikzpicture}
  \caption{Типовой цикл отладки.}
  \label{fig:debug-cycle}
\end{figure}

\begin{enumerate}
  \item \textbf{Воспроизвести} стабильно --- без
    воспроизведения отладка невозможна.
  \item \textbf{Изолировать}: бинарный поиск по коду
    (убрать половину), \cmd{printf} или \cmd{assert} в
    подозрительных местах.
  \item \textbf{Наблюдать} через gdb (значения) и
    Valgrind (память).
  \item \textbf{Починить} и \textbf{написать тест},
    чтобы не повторилось (\nameref{ch:15-testing-optimization}).
\end{enumerate}

\section{Типичные ошибки в самом отладчике}
\begin{itemize}
  \item Запустили Valgrind на оптимизированной сборке ---
    отчёт «гуляет»; пересоберите с \cmd{-O0 -g}.
  \item Не поставили \cmd{-g} --- в gdb видно только
    адреса, а не строки.
  \item Ищут баг, который невозможно воспроизвести ---
    сначала стабилизировать условия (фиксированные
    данные, seed для генератора случайных чисел).
\end{itemize}

\section{Связь с остальным курсом}

Отладка опирается на понимание памяти
(\nameref{ch:16-virtual-memory}) и процессов
(\nameref{ch:12-processes}). Valgrind в Docker
(\nameref{ch:19-docker-basics}) требует
\fpath{/proc} и работает медленнее, но применим;
тесты после фикса --- \nameref{ch:15-testing-optimization}.

\section{Задачи для самопроверки}
\begin{enumerate}
  \item Какие флаги \cmd{gcc} нужны для отладочной
    сборки?
  \item Чем \cmd{step} отличается от \cmd{next} в gdb?
  \item Что такое «definitely lost» в отчёте Valgrind?
  \item Приведите пример off-by-one ошибки в C.
  \item Почему Valgrind медленнее обычного запуска и
    когда это допустимо?
\end{enumerate}


\part{Взаимодействие процессов и виртуализация}
% Глава: 18-ipc
% Статус: наполнено (первый прогон).
\chapter{Взаимодействие процессов (IPC)}
\label{ch:18-ipc}

\section{Зачем процессам общаться}

Каждый процесс изолирован: своя память, свои файловые
дескрипторы. Но реальные системы --- это множество
процессов, которым нужно обмениваться данными. Способы
такого обмена называют \textbf{IPC} (Inter-Process
Communication).

Все сценарии IPC выполняются на вашей ВМ на \fpath{pglab.nsu.ru}:
общаются процессы одной вашей системы.

\section{Файлы: самый простой способ}

Процесс пишет в файл, другой читает. Медленно, но
просто и наглядно. Для медленных потоков данных
(конфигурации, отчёты) этого достаточно. Ограничение:
файловая система не «сигнализирует» о новых данных ---
второму процессу нужно опрашивать.

\section{Пайпы (pipes)}

\textbf{Пайп} --- односторонний канал.
\textbf{Анонимный пайп} живёт в памяти и существует
только между связанными процессами (родитель и
потомком):
\begin{minted}[bgcolor=gray!6]{bash}
# оболочка построит этот конвейер за вас:
ls -l | grep txt | wc -l
\end{minted}

\textbf{Именованный пайп} (FIFO) --- файл-устройство
на диске, доступное любому процессу:
\begin{minted}[bgcolor=gray!6]{bash}
mkfifo myfifo
# в одном терминале:
cat > myfifo
# в другом терминале:
echo hello > myfifo
\end{minted}

Пайпы работают как буфер фиксированного размера
(обычно 64 КБ): если буфер полон, пишущий блокируется.

\section{Сокеты}

\textbf{Сокет} --- обобщение пайпа: двунаправленный
канал между двумя точками. Бывают:
\begin{itemize}
  \item \textbf{Unix domain} (межпроцессные, в пределах
    одной системы) --- быстрые;
  \item \textbf{TCP/IP} (между машинами по сети) ---
    так работают веб-сервер и клиент, SSH,
    большинство клиент-серверных приложений.
\end{itemize}

\section{Общая память (shared memory)}

Два процесса могут «пришить» одну и ту же физическую
страницу к своим адресным пространствам
(\nameref{ch:16-virtual-memory}) и обмениваться
данными без копирования. Самый быстрый способ IPC, но
требует синхронизации (семафоры, мьютексы), чтобы
процессы не записывали одновременно.

\section{Очереди сообщений}

Готовый механизм ядра: процесс кладёт сообщение
в очередь, другой --- забирает. Проще, чем общая
память (ядро следит за целостностью), медленнее.

\section{Сигналы}

Особый, асинхронный способ (\nameref{ch:12-processes}):
короткое уведомление без передачи данных
(«остановись», «перезапусти»). Сигналы --- это
уведомление, а не канал данных.

\section{Выбор способа}

\begin{center}
\begin{tabular}{p{0.52\linewidth}l}
\toprule
\textbf{Ситуация} & \textbf{Способ} \\
\midrule
одноразовые данные & файл \\
быстрый поток между родителем и дочерним & анонимный пайп \\
связь между несвязанными процессами & именованный пайп или сокет \\
клиент-сервер в одной системе & Unix socket \\
клиент-сервер по сети & TCP/IP socket \\
максимальная скорость & shared memory + синхронизация \\
\bottomrule
\end{tabular}
\end{center}

\section{Связь с остальным курсом}

Анонимные пайпы собирает оболочка
(\nameref{ch:07-text-pipelines}); сокеты лежат в основе
SSH (\nameref{ch:02-linux-access}) и Docker API
(\nameref{ch:19-docker-basics}); ограничения на
ресурсы контейнеров --- cgroups
(\nameref{ch:23-docker-resources}).

\section{Задачи для самопроверки}
\begin{enumerate}
  \item Перечислите основные способы IPC и
    краткое назначение каждого.
  \item Чем анонимный пайп отличается от именованного?
  \item Когда общая память лучше сокета, а когда ---
    нет?
  \item Какие данные передаёт сигнал?
  \item Почему shared memory требует синхронизации?
\end{enumerate}

% Глава: 19-docker-basics
% Статус: наполнено (первый прогон).
\chapter{Docker: основы}
\label{ch:19-docker-basics}

\section{Контейнер и виртуальная машина}

\textbf{Виртуальная машина} виртуализует
\emph{оборудование}: внутри гостевая ОС с ядром,
драйверами, собственным пространством.
\textbf{Контейнер} изолирует \emph{процессы}: все
контейнеры на хосте делят одно ядро Linux, но имеют
отдельные файловые системы, сетевые
пространства и PID-пространства.

На курсе Docker-демон (хост) --- это ваша ВМ на сервере
\fpath{pglab.nsu.ru}: контейнеры изолируются внутри вашей системы,
разделяя её ядро и диск.

\begin{figure}[htb]
  \centering
  \begin{tikzpicture}[node distance=0.6cm]
    \node[layer, minimum width=6cm, minimum height=2.6cm] (vmhost) {};
    \node[font=\small, above=2pt of vmhost] {Виртуальные машины};
    \node[process, minimum width=2.4cm] (vm1) at (vmhost) {}; 
    \node[font=\small] at (vm1) {ВМ: гостевая ОС + ядро};
    \node[layer, below=1.6cm of vmhost, minimum width=6cm, minimum height=2.6cm] (cghost) {};
    \node[font=\small, above=2pt of cghost] {Контейнеры};
    \node[ext, minimum width=1.5cm] (c1) at ([yshift=-0.5cm]cghost.west) [xshift=0.7cm] {app};
    \node[ext, right=0.3cm of c1, minimum width=1.5cm] (c2) {app};
    \node[ext, right=0.3cm of c2, minimum width=1.5cm] (c3) {app};
    \node[font=\small, below=3pt of c2] {общее ядро хоста};
  \end{tikzpicture}
  \caption{ВМ изолирует целую ОС, контейнеры --- процессы поверх общего ядра.}
  \label{fig:vm-vs-ct}
\end{figure}

Преимущества контейнеров: быстрый запуск (секунды),
малый размер (сотни МБ), переносимость «от
разработчика до продакшена».

\section{Архитектура Docker}

\begin{itemize}
  \item \textbf{Docker Engine} (демон \cmd{dockerd}) ---
    управляющий процесс на хосте;
  \item \textbf{клиент} (\cmd{docker}) --- команда,
    которой вы пользуетесь; общается с демоном по
    API;
  \item \textbf{образ} (image) --- неизменяемый набор
    файлов + метаданные; \textbf{контейнер} ---
    работающий экземпляр образа.
\end{itemize}

\section{Ключевые команды}

\begin{minted}[bgcolor=gray!6]{bash}
docker run hello-world        # скачать и запустить тестовый образ
docker run -it ubuntu bash    # запустить интерактивный контейнер
docker ps                     # работающие контейнеры
docker ps -a                  # все контейнеры, включая остановленные
docker images                 # локальные образы
docker exec -it <id> bash     # выполнить команду в работающем контейнере
docker logs <id>              # stdout/stderr контейнера
docker stop <id>              # остановить контейнер
docker start <id>             # запустить остановленный контейнер
docker rm <id>                # удалить остановленный контейнер
docker rmi <image>            # удалить образ
\end{minted}

\cmd{docker run} с флагами:
\begin{minted}[bgcolor=gray!6]{bash}
docker run -d -p 8080:80 --name web nginx   # в фоне, маппинг портов
docker run --rm -it alpine sh               # удалить при выходе
docker run --name app myimage:1.0           # по имени и тегу образа
\end{minted}

\section{Образы и реестры}

\textbf{Реестр} (registry) --- хранилище образов;
основной публичный --- Docker Hub.
\textbf{Тег} (tag) --- метка версии в имени образа:
\cmd{nginx:1.25}, \cmd{nginx:latest}.
\begin{minted}[bgcolor=gray!6]{bash}
docker pull nginx:1.25           # скачать конкретную версию
docker tag myimage:1.0 me/myimage:1.0   # переименовать для push
docker push me/myimage:1.0       # загрузить, после login
\end{minted}

\section{Типичные ошибки}
\begin{itemize}
  \item \emph{permission denied} на \fpath{/var/run/docker.sock}
    --- пользователь не в группе \cmd{docker}; добавьте
    (\cmd{sudo usermod -aG docker \$USER}), перезайдите.
  \item Контейнер сразу завершился --- проверьте
    \cmd{docker logs}; часто причина --- отсутствующий
    \cmd{CMD} в Dockerfile.
  \item Порт «не отвечает» --- не задан \cmd{-p host:container}.
  \item Изменения в контейнере потерялись после
    перезапуска --- данные хранятся в слое контейнера;
    для сохранения нужны volumes (\nameref{ch:21-docker-data}).
\end{itemize}

\section{Связь с остальным курсом}

Создание образа --- Dockerfile
(\nameref{ch:20-dockerfile}); хранение данных ---
\nameref{ch:21-docker-data}; несколько контейнеров ---
\nameref{ch:22-docker-compose}; ограничения ресурсов ---
\nameref{ch:23-docker-resources}. Внутрь контейнера
можно войти \cmd{docker exec} и поработать
в tmux (\nameref{ch:10-tmux}) или Vim
(\nameref{ch:09-vim}).

\section{Задачи для самопроверки}
\begin{enumerate}
  \item Чем контейнер отличается от виртуальной
    машины?
  \item Что такое образ и что такое контейнер?
  \item Как запустить интерактивный контейнер с
    bash и как из него выйти?
  \item Как посмотреть логи завершившегося
    контейнера?
  \item Зачем нужны теги образов?
\end{enumerate}

% Глава: 20-dockerfile
% Статус: наполнено (первый прогон).
\chapter{Dockerfile}
\label{ch:20-dockerfile}

\section{Что такое Dockerfile}

\textbf{Dockerfile} --- текст с инструкциями, как
построить образ. Команда \cmd{docker build} выполняет
инструкции по порядку, каждая --- отдельный
\textbf{слой}.

Собираем образы на своей ВМ: \cmd{docker build} запускается на
\fpath{pglab.nsu.ru} и кэширует слои на её диске.

\begin{minted}[bgcolor=gray!6]{Dockerfile}
# базовый образ
FROM ubuntu:22.04

# метаданные
LABEL maintainer="dev@example.org"

# рабочий каталог внутри контейнера
WORKDIR /app

# установить инструменты; каждый RUN создаёт отдельный слой
RUN apt-get update && apt-get install -y --no-install-recommends gcc make
RUN rm -rf /var/lib/apt/lists/*

# скопировать файлы проекта
COPY src/ ./src/

# собрать программу
RUN make

# точка входа
CMD ["./app"]
\end{minted}

\section{Основные инструкции}

\begin{center}
\begin{tabular}{lp{0.62\linewidth}}
\toprule
\textbf{Инструкция} & \textbf{Назначение} \\
\midrule
\cmd{FROM}     & базовый образ; начинается с него сборка \\
\cmd{WORKDIR}  & каталог по умолчанию внутри контейнера \\
\cmd{RUN}      & выполнить команду при сборке (создаёт слой) \\
\cmd{COPY}     & скопировать файлы из контекста сборки \\
\cmd{ADD}      & как COPY + распаковка архивов и загрузка по URL \\
\cmd{ENV}      & задать переменную окружения в образе \\
\cmd{ARG}      & параметр, передаваемый при \cmd{docker build} \\
\cmd{LABEL}    & метаданные образа \\
\cmd{EXPOSE}   & документация о портах (не пробрасывает!) \\
\cmd{ENTRYPOINT} & фиксированная команда контейнера \\
\cmd{CMD}      & аргументы по умолчанию для ENTRYPOINT \\
\cmd{USER}     & пользователь, от которого работает контейнер \\
\bottomrule
\end{tabular}
\end{center}

\textbf{CMD} vs \textbf{ENTRYPOINT}: \cmd{CMD} легко
переопределить аргументами запуска; \cmd{ENTRYPOINT} ---
``главная'' команда, аргументы передаются ей.

\section{Слои и кэш}

Каждая инструкция создаёт слой. Docker кэширует
слои: если строки выше не менялись, их слои берутся
из кэша, и сборка идёт быстро. Поэтому
\textbf{порядок инструкций важен}:
\begin{itemize}
  \item сначала то, что редко меняется (зависимости);
  \item потом то, что часто (исходный код);
  \item так при изменении кода пересобирается только
    верхний слой.
\end{itemize}

\section{\cmd{COPY} и \cmd{ADD}}

\cmd{COPY} --- скопировать файл из каталога сборки.
\cmd{ADD} делает то же, но дополнительно распаковывает
архивы и может загрузить файл по URL. Для обычного
копаирования файлов проекта используйте \cmd{COPY}.

\section{Многоступенчатая сборка (multi-stage)}

Один Dockerfile может содержать несколько
\cmd{FROM}:
\begin{minted}[bgcolor=gray!6]{Dockerfile}
FROM golang:1.21 AS build
WORKDIR /src
COPY . .
RUN go build -o app .

FROM alpine:3.18
COPY --from=build /src/app /app
ENTRYPOINT ["/app"]
\end{minted}

Первый слой --- «сборочный» (с компилятором), второй ---
лёгкий, только с готовым бинарником. Размер
конечного образа резко меньше.

\section{\cmd{.dockerignore}}

Файл \fpath{.dockerignore} рядом с Dockerfile
определяет, что \emph{не} копировать в контекст
сборки (.git, локальные логи, артефакты).
Уменьшает размер и время сборки.

\section{Типичные ошибки}
\begin{itemize}
  \item \cmd{RUN apt-get update} и \cmd{apt-get install}
    в разных \cmd{RUN} --- кэш не поможет; объедините
    в одну инструкцию;
  \item Забыли \cmd{WORKDIR} --- файлы копировались в
    \fpath{/} вместо нужного каталога;
  \item Образ «огромный» --- кэш пакетного менеджера не
    чистится после установки (см.\ пример с \cmd{rm -rf}
    выше), либо не используется multi-stage;
  \item \cmd{EXPOSE} «не работает» --- он не пробрасует
    порт; нужен \cmd{-p} при \cmd{docker run}.
\end{itemize}

\section{Связь с остальным курсом}

Dockerfile --- описание образа; хранение данных ---
\nameref{ch:21-docker-data}; несколько контейнеров ---
\nameref{ch:22-docker-compose}; сборка C-проекта внутри
образа --- \nameref{ch:13-c-compilation} и
\nameref{ch:14-make}.

\section{Задачи для самопроверки}
\begin{enumerate}
  \item Чем \cmd{COPY} отличается от \cmd{ADD}?
  \item Что такое слой и почему порядок инструкций
    важен?
  \item Чем \cmd{CMD} отличается от \cmd{ENTRYPOINT}?
  \item Как сделать образ «меньше», используя
    multi-stage?
  \item Зачем нужен \cmd{.dockerignore}?
\end{enumerate}

% Глава: 21-docker-data
% Статус: наполнено (перый прогон).
\chapter{Данные в Docker: тома и монтирование}
\label{ch:21-docker-data}

\section{Проблема: данные в контейнере теряются}

Всё, что записано в файловую систему контейнера
(не в том, что было в образе), живёт в \emph{слое
контейнера}. Когда контейнер удаляется --- данные
исчезают. Для сохранения нужны \textbf{тома}
(volumes) и \textbf{монтирование} (mounts).

Тома размещаются на диске вашей ВМ на \fpath{pglab.nsu.ru};
учитывайте, что там 20--40~ГБ, и следите за местом.

\section{Три способа хранения}

\begin{center}
\begin{tabular}{lp{0.24\linewidth}lp{0.24\linewidth}}
\toprule
\textbf{Способ} & \textbf{Хранение} & \textbf{Управление} & \textbf{Кто управляет} \\
\midrule
named volume & в хранилище Docker & Docker & Docker \\
bind mount   & произвольный каталог хоста & вы & вы \\
tmpfs        & в RAM хоста & Docker & Docker \\
\bottomrule
\end{tabular}
\end{center}

\section{Именованные тома (named volumes)}

Создаются и управляются Docker:
\begin{minted}[bgcolor=gray!6]{bash}
docker volume create mydata          # создать
docker volume ls                     # список
docker volume inspect mydata         # подробно
docker volume rm mydata              # удалить, только пустой
\end{minted}

Подключить к контейнеру:
\begin{minted}[bgcolor=gray!6]{bash}
docker run -d -v mydata:/var/lib/app-data myimage
\end{minted}

\cmd{-v host:container} --- общий синтаксис:
что хост-путь (или имя тома), что путь внутри
контейнера.

\section{Bind mounts}

Монтирование \emph{произвольного} каталога хоста:
\begin{minted}[bgcolor=gray!6]{bash}
docker run -it -v /home/me/project:/app ubuntu ls /app
\end{minted}

Полезно для разработки: меняете файл на хосте ---
он сразу виден в контейнере.

\section{tmpfs}

\textbf{tmpfs} --- каталог в RAM; исчезает при
остановке контейнера:
\begin{minted}[bgcolor=gray!6]{bash}
docker run -it --tmpfs /tmp:size=100m ubuntu ls -la /tmp
\end{minted}

\section{Схема: хост, образ, контейнер, том}

\begin{figure}[htb]
  \centering
  \begin{tikzpicture}[node distance=1.4cm]
    \node[storage, minimum width=4cm] (vol) {named volume\\[0.2em] (on host disk)};
    \node[process, above right=1.2cm and 1cm of vol] (ctr) {container\\[0.2em] (ephemeral)};
    \node[data, above left=1.2cm and 1cm of vol] (img) {image (read-only)};
    \draw[flow, <->] (vol) -- node[right, font=\small]{mount} (ctr);
    \draw[flow, ->] (img) -- node[left, font=\small]{copy up} (ctr);
  \end{tikzpicture}
  \caption{Образ --- только чтение, слой контейнера --- временный,
    том --- постоянное хранилище.}
  \label{fig:docker-data}
\end{figure}

\section{Практика: база данных в Docker}

Типовой сценарий: образ с приложением + отдельный том
для данных:
\begin{minted}[bgcolor=gray!6]{bash}
docker volume create pg_data
docker run -d \\
  --name pg \\
  -v pg_data:/var/lib/postgresql/data \\
  -e POSTGRES_PASSWORD=secret \\
  postgres:16
\end{minted}

Даже если контейнер удалить, данные останутся в томе
\cmd{pg_data}.

\section{Типичные ошибки}
\begin{itemize}
  \item Данные «пропали» после \cmd{docker rm} --- не
    подключён том;
  \item Bind mount на каталог, которого нет на хосте ---
    Docker создаст его \emph{пустым};
  \item Права: приложение в контейнере работает от
    \cmd{root}, а на хосте каталог принадлежит вам ---
    возможны конфликты прав.
\end{itemize}

\section{Связь с остальным курсом}

Тома используются в \cmd{docker-compose}
(\nameref{ch:22-docker-compose}); ограничения на
ресурсы контейнеров --- \nameref{ch:23-docker-resources};
файловая система хоста --- \nameref{ch:04-filesystem}.

\section{Задачи для самопроверки}
\begin{enumerate}
  \item Почему данные в контейнере теряются после
    его удаления?
  \item Чем named volume отличается от bind mount?
  \item Как подключить том к контейнеру?
  \item Что такое tmpfs и когда его используют?
  \item Как сохранить данные PostgreSQL при удалении
    контейнера?
\end{enumerate}

% Глава: 22-docker-compose
% Статус: наполнено (первый прогон).
\chapter{Docker Compose}
\label{ch:22-docker-compose}

\section{Зачем Compose}

Когда приложение --- это несколько контейнеров
(веб, база данных, очередь), описывать их
отдельными \cmd{docker run} неудобно.
\textbf{Docker Compose} описывает всю систему в одном
YAML-файле и запускает её одной командой.

Оркеструем такие системы на своей ВМ: Compose поднимает
все контейнеры на \fpath{pglab.nsu.ru} одной командой.

\section{Анатомия \fpath{docker-compose.yml}}

\begin{minted}[bgcolor=gray!6]{bash}
services:
  web:
    build: .
    ports:
      - "8080:80"
    depends_on:
      - db
    environment:
      - DATABASE_URL=postgres://db:5432
  db:
    image: postgres:16
    volumes:
      - dbdata:/var/lib/postgresql/data
    environment:
      POSTGRES_PASSWORD: secret

volumes:
  dbdata:
\end{minted}

Ключевые поля:
\begin{center}
\begin{tabular}{lp{0.62\linewidth}}
\toprule
\textbf{Поле} & \textbf{Назначение} \\
\midrule
\cmd{services}      & словарь сервисов (контейнеров) \\
\cmd{image}         & готовый образ \\
\cmd{build}         & собрать из Dockerfile (путь) \\
\cmd{ports}         & проброс портов (host:container) \\
\cmd{volumes}       & тома (именованные или bind mount) \\
\cmd{environment}   & переменные окружения \\
\cmd{depends\_on}   & порядок запуска (запустит сначала) \\
\cmd{networks}      & пользовательские сети \\
\cmd{command}       & переопределить CMD образа \\
\bottomrule
\end{tabular}
\end{center}

\section{Основные команды}

\begin{minted}[bgcolor=gray!6]{bash}
docker compose up            # собрать при необходимости и запустить все сервисы
docker compose up -d         # запустить в фоне
docker compose ps            # список работающих сервисов
docker compose logs web      # логи сервиса "web"
docker compose exec web bash # открыть оболочку в сервисе "web"
docker compose stop          # остановить все
docker compose down          # остановить и удалить контейнеры
docker compose down -v       # также удалить тома, осторожно!
\end{minted}

\section{Сети и DNS}

Компоненты в одном compose-файле автоматически
попадают в одну сеть; имя сервиса --- его DNS-имя:
веб обращается к базе как \cmd{db} (не \cmd{localhost}).

\section{Типичные ошибки}
\begin{itemize}
  \item \cmd{depends\_on} --- только про
    \emph{запуск}, не про \emph{готовность}: база может
    ещё не быть готова принимать подключения, когда
    стартует веб; используйте \cmd{healthcheck} или
    retry в коде;
  \item Порт занят на хосте --- \cmd{docker compose ps}
    покажет; смените host-порт в \cmd{ports};
  \item \cmd{docker compose down -v} удаляет тома и
    данные --- используйте с осторожностью
    (\nameref{ch:21-docker-data});
  \item Несовместимые версии \cmd{compose} и
    \cmd{docker} --- обновите инструмент.
\end{itemize}

\section{Связь с остальным курсом}

Compose комбинирует Dockerfile
(\nameref{ch:20-dockerfile}) и тома
(\nameref{ch:21-docker-data}); ограничения ресурсов
для каждого сервиса --- \nameref{ch:23-docker-resources};
запуск нескольких контейнеров на хосте ---
\nameref{ch:19-docker-basics}.

\section{Задачи для самопроверки}
\begin{enumerate}
  \item Чем Compose удобнее, чем несколько
    \cmd{docker run}?
  \item Какие ключевые поля у сервиса в
    \fpath{docker-compose.yml}?
  \item Как запустить систему в фоне и посмотреть
    логи одного сервиса?
  \item Чем \cmd{depends\_on} не гарантирует
    (и что нужно добавить)?
  \item Как сервисам обращаться друг к другу по
    имени?
\end{enumerate}

% Глава: 23-docker-resources
% Статус: наполнено (первый прогон).
\chapter{Ресурсы контейнеров}
\label{ch:23-docker-resources}

\section{Изоляция ресурсов: cgroups}

Контейнеры делят ядро хоста; чтобы один «не съел»
всё, Linux использует \textbf{cgroups} (control
groups) --- механизм ядра, ограничивающий ресурсы
(память, CPU, IO) для группы процессов. Docker
управляет cgroups за вас через флаги.

Лимиты задают долю ресурсов вашей ВМ на \fpath{pglab.nsu.ru},
чтобы один контейнер не забрал всё у соседних.

\section{Лимиты CPU}

\begin{minted}[bgcolor=gray!6]{bash}
docker run --cpus 1.5 myimage          # ограничить 1.5 ядрами
docker run --cpuset-cpus 0,1 myimage   # закрепить за конкретными CPU
\end{minted}

\cmd{--cpus} задаёт долю (дробное число допустимо),
\cmd{--cpuset-cpus} --- список конкретных ядер.

\section{Лимиты памяти}

\begin{minted}[bgcolor=gray!6]{bash}
docker run --memory 512m myimage          # максимум 512 МиБ
docker run --memory 512m --memory-swap 1g myimage  # + подкачка (swap)
\end{minted}

Если процесс превысит лимит памяти, ядро сработает
\textbf{OOM-killer} и завершит процесс с ошибкой
(\cmd{Killed}). Проверьте логи:
\cmd{docker inspect} покажет \cmd{OOMKilled: true}.

\section{Мониторинг}

\begin{minted}[bgcolor=gray!6]{bash}
docker stats             # живая статистика CPU и памяти по контейнерам
docker stats --no-stream # однократно
docker top <id>          # процессы внутри контейнера
docker inspect <id>      # полная конфигурация и лимиты
\end{minted}

\section{Типовая схема ограничений}

\begin{figure}[htb]
  \centering
  \begin{tikzpicture}[node distance=1.3cm]
    \node[layer, minimum width=7cm, minimum height=2.8cm] (host) {};
    \node[font=\small] at (host) {хост: CPU + RAM};
    \node[process, minimum width=1.6cm] (c1) at ([yshift=-0.4cm]host.west) [xshift=0.9cm] {app 1};
    \node[process, right=0.3cm of c1, minimum width=1.6cm] (c2) {app 2};
    \node[font=\tiny, below=1pt of c1] {cpus=1.0, mem=512m};
    \node[font=\tiny, below=1pt of c2] {cpus=0.5, mem=256m};
  \end{tikzpicture}
  \caption{Каждый контейнер имеет лимиты CPU и памяти (cgroups).}
  \label{fig:cgroups}
\end{figure}

\section{Типичные ошибки}
\begin{itemize}
  \item Контейнер убивается «без причины» --- смотрите
    \cmd{OOMKilled}; увеличьте \cmd{--memory} или
    оптимизируйте память приложения;
  \item Лимит CPU слишком низкий --- приложение
    «замедляется»; измерьте \cmd{docker stats};
  \item Не заданы лимиты --- один контейнер может
    занять всю память хоста;
  \item \cmd{--memory-swap} без \cmd{--memory} ---
    некорректная конфигурация.
\end{itemize}

\section{Связь с остальным курсом}

Лимиты настраиваются и в \cmd{docker-compose}
(\nameref{ch:22-docker-compose}); мониторинг процессов
--- \nameref{ch:12-processes}; распределение памяти ---
\nameref{ch:16-virtual-memory}.

\section{Задачи для самопроверки}
\begin{enumerate}
  \item Что такое cgroups и зачем они нужны
    контейнерам?
  \item Как ограничить контейнер 1 CPU и 512 МБ
    памяти?
  \item Как узнать, был ли контейнер убит
    OOM-killer'ом?
  \item Чем \cmd{--cpus} отличается от
    \cmd{--cpuset-cpus}?
  \item Как посмотреть использование ресурсов
    контейнеров в реальном времени?
\end{enumerate}


\part{Среда разработчика}
% Глава: 24-security
% Статус: задел (кратко), будет дополнено.
\chapter{Основы безопасности}
\label{ch:24-security}

\emph{Раздел будет дополнен в следующих версиях
пособия. Пока --- краткая вводка и ключевые темы.}

Безопасность здесь --- прежде всего безопасность доступа к
серверу лаборатории \fpath{pglab.nsu.ru} и к вашей ВМ на нём.

\section{Модель пользователя и привилегий}

Каждый процесс работает от имени пользователя
(\nameref{ch:05-permissions}); \cmd{sudo} даёт
временные права суперпользователя. Правила:
\begin{itemize}
  \item минимальные привилегии: запускаться под
    тем пользователем, который действительно
    нужен;
  \item не хранить секреты в коде и в
    \fpath{.gitignore}-не-защитённых местах;
  \item SSH-ключи предпочтительнее паролей
    (\nameref{ch:02-linux-access}).
\end{itemize}

\section{Что будет в полном разделе}
\begin{itemize}
  \item Файловые права и SELinux/AppArmor
    на практике;
  \item Защита SSH: отключение входа по паролю,
    AllowUsers, fail2ban;
  \item Работа с секретами (переменные окружения,
    файловые ключи, Vault-подход);
  \item Сетевая изоляция: firewall (nftables/
    ufw), группы безопасности;
  \item Обновления и патчи, управление зависимостями.
\end{itemize}

\section{Задачи для самопроверки}
\begin{enumerate}
  \item Почему SSH-ключи безопаснее паролей?
  \item Какие минимальные привилегии должны быть
    у сервиса в контейнере?
  \item Где \emph{не} следует хранить секреты?
\end{enumerate}

% Глава: 25-git
% Статус: наполнено (первый прогон).
\chapter{Git: системы контроля версий}
\label{ch:25-git}

\section{Зачем нужен контроль версий}

Версия кода --- это состояние проекта в какой-то
момент. Без контроля версий трудно: вернуть старый
вариант, понять, кто и что менял, работать над
несколькими задачами параллельно, делиться кодом.
\textbf{Git} --- распределённая система контроля
версий (VCS): полная история на каждой машине,
быстрые локальные операции, ветвление как базовый
механизм.

Работаем с Git на своей ВМ на \fpath{pglab.nsu.ru}: репозиторий
живёт на сервере, а удалённым хостом служит, например, GitHub
или внутренний Git на сети университета.

\section{Распределённая модель}

В централизованных системах (SVN) есть один центральный
сервер. В Git \emph{каждая копия} --- полноценное
хранилище со всей историей; ``сервер''
(\cmd{github.com}, \cmd{gitlab.com}) нужен для
координации, а не как обязательная точка. Работать
можно офлайн, коммиты локальные, а \cmd{push} ---
это синхронизация.

\section{Три области}

\begin{figure}[htb]
  \centering
  \begin{tikzpicture}[node distance=0.9cm]
    \node[data, minimum width=2.9cm] (wd) {working tree\\[0.2em] (files on disk)};
    \node[process, right=of wd, minimum width=2.9cm] (idx) {index\\[0.2em] (staging area)};
    \node[process, right=of idx, minimum width=2.9cm] (repo) {repository\\[0.2em] (commits)};
    \node[process, right=of repo, minimum width=2.9cm] (remote) {remote\\[0.2em] (origin)};
    \draw[flow] (wd) -- node[above, font=\small]{git add} (idx);
    \draw[flow] (idx) -- node[above, font=\small]{git commit} (repo);
    \draw[flow, <->] (repo) -- node[above, font=\small]{push / pull} (remote);
  \end{tikzpicture}
  \caption{Три области Git: рабочая директория, индекс, хранилище.}
  \label{fig:git-areas}
\end{figure}

\begin{itemize}
  \item \textbf{working tree} --- файлы на диске;
  \item \textbf{index} (staging area) --- набор
    изменений, которые попадут в следующий коммит;
  \item \textbf{repository} --- цепочка коммитов;
  \item \textbf{remote} --- удалённое хранилище
    (\cmd{origin}).
\end{itemize}

\section{Первые шаги}

\begin{minted}[bgcolor=gray!6]{bash}
git init                 # создать здесь новый репозиторий
git clone <url>          # скопировать существующий репозиторий
cd myrepo
\end{minted}

\section{Базовый цикл}

\begin{minted}[bgcolor=gray!6]{bash}
git status               # что изменилось
git add file.txt         # добавить один файл в индекс
git add .                # добавить все изменения этого каталога
git commit -m "message"  # создать коммит
git log --oneline        # краткая история
\end{minted}

\textbf{Коммит} --- атомарный снимок: кто, когда, что
изменил и почему. Текст коммита описывает \emph{что}
и \emph{зачем}.

\section{Ветвление и слияние}

\begin{minted}[bgcolor=gray!6]{bash}
git branch feature-x     # создать ветку
git switch feature-x     # переключиться на неё
git switch -c feature-y  # создать и переключиться за один шаг
git merge feature-x      # слить ветку в текущую
git branch -d feature-x  # удалить после слияния
\end{minted}

\textbf{merge} --- объединяет историю двух веток,
создавая коммит-слияние. \textbf{rebase} ---
``пересаживает'' вашу ветку на другую, история
линейнее, но коммиты переписываются.

\begin{minted}[bgcolor=gray!6]{bash}
git rebase main          # переиграть эту ветку поверх main
\end{minted}

\begin{tipbox}
Правило: \cmd{merge} для публичных веток (честная
история), \cmd{rebase} для личных веток до
публикации (чистая история). Никогда не
переписывайте уже опубликованную историю.
\end{tipbox}

\section{Работа с удалённым хранилищем}

\begin{minted}[bgcolor=gray!6]{bash}
git remote add origin <url>   # указать удалённый репозиторий
git push -u origin main       # отправить и задать upstream
git pull                      # забрать и слить с upstream
git fetch                     # только скачать, без слияния
git branch -r                 # список удалённых веток
\end{minted}

\section{\cmd{.gitignore}}

Файл \fpath{.gitignore} --- какие файлы Git
\emph{не} отслеживает:
\begin{minted}[bgcolor=gray!6]{bash}
# артефакты сборки
*.o
*.a
out/
build/

# редактор и IDE
.vscode/
.idea/
*.swp

# секреты и локальная конфигурация
.env
config.local.yml
\end{minted}

\section{Просмотр изменений и история}

\begin{minted}[bgcolor=gray!6]{bash}
git diff                 # недобавленные изменения
git diff --staged        # добавленные в индекс
git diff HEAD            # все изменения относительно последнего коммита
git log --oneline --graph --all   # наглядная история
git show <commit>        # что изменил коммит
\end{minted}

\section{Откат и исправление}

\begin{minted}[bgcolor=gray!6]{bash}
git reset --soft HEAD~1    # отменить последний коммит, оставить изменения в индексе
git reset --mixed HEAD~1   # отменить последний коммит, не добавляя в индекс
git reset --hard HEAD~1    # отменить последний коммит и выбросить изменения, осторожно!
git revert <commit>        # создать новый коммит, отменяющий указанный
\end{minted}

\cmd{reset} двигает указатель ветки; \cmd{revert} ---
безопасная операция, создающая новый коммит с
обратным эффектом. Для публичной истории используйте
\cmd{revert}.

\section{\cmd{.gitconfig} и идентификация}

\begin{minted}[bgcolor=gray!6]{bash}
git config --global user.name "Jane Doe"
git config --global user.email "jane@example.org"
git config --global init.defaultBranch main
git config --global core.editor vim
\end{minted}

\section{Типичные ошибки}
\begin{itemize}
  \item Забыли \cmd{git add} перед \cmd{commit} ---
    коммит пустой;
  \item В коммите ``fixed bug'' без объяснения ---
    невозможно понять, \emph{какой} баг и \emph{почему}
    так;
  \item \cmd{git push --force} на общую ветку ---
    история переписывается, у коллег конфликты;
  \item Секреты (пароли, токены) в репозитории ---
    добавьте в \cmd{.gitignore} и \emph{измените
    секрет};
  \item \cmd{git merge} и \cmd{git rebase} в одном
    рабочем процессе --- путаница; выберите один
    подход.
\end{itemize}

\section{Связь с остальным курсом}

Git --- это то, как делятся проекты в команде
(\nameref{ch:27-scripting-ide}) и что
интегрируют в CI (\nameref{ch:26-ci}). В Docker
(\nameref{ch:19-docker-basics}) образы часто
тегируются номером коммита.

\section{Задачи для самопроверки}
\begin{enumerate}
  \item Что такое коммит и чем он отличается от
    снимка в резервном копировании?
  \item Чем \cmd{merge} отличается от \cmd{rebase}?
    Когда что использовать?
  \item Как создать ветку, переключиться в неё и
    слить обратно?
  \item Что делает \cmd{git reset --hard} и почему он
    опасен?
  \item Какие файлы обычно попадают в
    \cmd{.gitignore}?
  \item Чем \cmd{pull} отличается от \cmd{fetch}?
\end{enumerate}

% Глава: 26-ci
% Статус: задел (кратко), будет дополнено.
\chapter{CI/CD: автоматизация сборки и поставки}
\label{ch:26-ci}

\emph{Раздел будет дополнен в следующих версиях
пособия. Пока --- краткая вводка и ключевые темы.}

\section{Что такое CI/CD}

\textbf{CI} (Continuous Integration) --- автоматический
запуск сборки и тестов на каждый push в
репозиторий. \textbf{CD} (Continuous
Delivery/Deployment) --- автоматическое
подготовка/выкатка релиза. Инструменты: GitHub
Actions, GitLab CI, Jenkins.

Обычно CI работает в облаке (GitHub Actions), но
самохостный runner можно развернуть и на сервере
лаборатории, чтобы сборки проходили в той же среде, что и
ваши ВМ на \fpath{pglab.nsu.ru}.

\section{Типовой пайплайн}

\begin{minted}[bgcolor=gray!6]{bash}
# псевдоструктура CI-конвейера
on: push
jobs:
  build:
    - checkout code
    - build project (make / docker build)
    - run tests
    - (optional) deploy
\end{minted}

\section{Что будет в полном разделе}
\begin{itemize}
  \item Конфигурация GitHub Actions / GitLab CI;
  \item Матрица сборки (несколько версий ОС/зависимостей);
  \item Кэширование зависимостей;
  \item Публикация артефактов и образов;
  \item Деплой в staging и прод.
\end{itemize}

\section{Связь с остальным курсом}

CI опирается на \cmd{make} (\nameref{ch:14-make}),
тестирование
(\nameref{ch:15-testing-optimization}), Docker
(\nameref{ch:19-docker-basics}) и Git
(\nameref{ch:25-git}).

\section{Задачи для самопроверки}
\begin{enumerate}
  \item Что даёт CI команде разработчиков?
  \item Какие этапы обычно есть в пайплайне?
  \item Какую роль играет Docker в CD?
\end{enumerate}

% Глава: 27-scripting-ide
% Статус: задел (кратко), будет дополнено.
\chapter{Скрипты оболочки и IDE}
\label{ch:27-scripting-ide}

\emph{Раздел будет дополнен в следующих версиях
пособия. Пока --- краткая вводка и ключевые темы.}

Скрипты и IDE мы используем на своей ВМ на \fpath{pglab.nsu.ru}:
собираем и запускаем проекты там, где они и живут.

\section{Скрипты оболочки}

Повторяющиеся команды логично собирать в скрипт
\fpath{run.sh}:
\begin{minted}[bgcolor=gray!6]{bash}
#!/usr/bin/env bash
# сборка и тестирование проекта
set -euo pipefail

make clean
make
make test
\end{minted}

\cmd{set -e} --- остановка при ошибке,
\cmd{set -u} --- ошибка на несуществующей
переменной, \cmd{set -o pipefail} --- ошибка в
пайпе, если любая команда пайпа упала.

\section{IDE для C и системной разработки}

\begin{itemize}
  \item \textbf{VS Code} + расширения (C/C++,
    clangd, CMake) --- популярен, настраиваем;
  \item \textbf{CLion} --- коммерческая IDE,
    мощный рефакторинг;
  \item \textbf{Neovim} + LSP --- для тех, кто
    хочет работать в терминале
    (\nameref{ch:09-vim}).
\end{itemize}

\section{Что будет в полном разделе}
\begin{itemize}
  \item Синтаксис bash: переменные, циклы,
    функции, условия;
  \item Аргументы командной строки (\cmd{\$1},
    \cmd{getopts});
  \item Обработка ошибок и логирование;
  \item Настройка IDE под C/C++: отладчик,
    компилятор, форматирование;
  \item LSP и автодополнение.
\end{itemize}

\section{Связь с остальным курсом}

Скрипты используют конвейеры
(\nameref{ch:07-text-pipelines}) и оболочки
(\nameref{ch:03-shell}); запуск в Docker
(\nameref{ch:19-docker-basics}); история изменений ---
Git (\nameref{ch:25-git}).

\section{Задачи для самопроверки}
\begin{enumerate}
  \item Зачем \cmd{set -euo pipefail} в скрипте?
  \item Как передать аргументы в bash-скрипт и
    как к ним обратиться?
  \item Чем IDE отличается от редактора текста?
\end{enumerate}


% ---- Приложения: отдельная часть (не внутри «Среды разработчика») ----
\part{Приложения}
\appendix
% Приложение: Глоссарий
\chapter{Глоссарий}
\label{app:glossary}

% Термины, используемые в методичке. Заполняется по мере написания глав.
% Формат: \item \textbf{Термин} --- значение.
\begin{description}
  \item{bash} --- стандартная командная оболочка Unix/Linux.
  \item{stdin/stdout/stderr} --- стандартные потоки ввода, вывода и ошибок.
  \item{пайп (pipe)} --- механизм передачи вывода одной команды во вход другой (\texttt{|}).
  \item{Vim} --- полноэкранный текстовый редактор.
  \item{tmux} --- мультплексер терминала.
  \item{Docker} --- платформа контейнеризации.
  % TODO: пополнять по мере написания глав.
\end{description}

% Приложение: Шпаргалки
\chapter{Шпаргалки по командам}
\label{app:cheatsheets}

% Сводные таблицы команд. Заполняется по мере написания глав.
% Каждая таблица — отдельная section, связанная с главой через \ref.

\section{Файловая система и навигация}
% \ref{ch:04-filesystem}
\begin{center}
\begin{tabular}{@{}ll@{}}
\texttt{ls}      & список содержимого каталога \\
\texttt{cd}      & смена каталога \\
\texttt{pwd}     & показать текущий каталог \\
\texttt{mkdir}   & создать каталог \\
\texttt{cp/mv/rm} & копирование / перемещение / удаление \\
\texttt{find}    & поиск файлов \\
\end{tabular}
\end{center}

\section{Текстовые конвейеры}
% \ref{ch:07-text-pipelines}
\begin{center}
\begin{tabular}{@{}ll@{}}
\texttt{grep} & поиск по регулярным выражениям \\
\texttt{awk}  & обработка по полям \\
\texttt{sed}  & преобразование текста \\
\texttt{sort} & сортировка \\
\texttt{head/tail/less} & просмотр файлов \\
\end{tabular}
\end{center}

% TODO: добавлять разделы по мере написания глав.

% Приложение: Список литературы (формат ГОСТ)
\chapter{Список литературы}
\label{app:references}

% Основной список оформлен по ГОСТ в виде thebibliography.
% По мере написания глав пополняйте; нумерация сквозная.
\begin{thebibliography}{99}

\bibitem{tannenbaum}
Таненбаум, Э.~С. Архитектура компьютера : [пер. с англ.] / Э.~Таненбаум, Т.~Остин. --- 6-е изд. --- Санкт-Петербург [и др.] : ПИТЕР, 2014. --- 811 с. : ил. ; 24 см. --- (Классика Computer Science). --- ISBN 978-5-496-00337-7.

\bibitem{krahotkina}
Архитектура ЭВМ : учебное пособие / авт.-сост. Е.~В. Крахоткина, В.~И. Терехин ; Северо-Кавказский федеральный университет. --- Ставрополь : СКФУ, 2015. --- 80 с. --- Режим доступа: по подписке. --- URL: \url{https://biblioclub.ru/index.php?page=book\&id=457862}. --- Текст : электронный.

\bibitem{dikov}
Диков, А.~В. Компьютер изнутри : учебное пособие : [16+] / А.~В. Диков. --- Москва ; Берлин : Директ-Медиа, 2015. --- 127 с. : ил., схем. --- ISBN 978-5-4475-5530-6. --- Текст : электронный.

% TODO: пополнять по мере написания глав (shell, Docker, Git, C/компиляция).
\end{thebibliography}

\section{Интернет-ресурсы}
\begin{description}
  \item{Страница курса} --- учебные материалы и задания на Google Classroom
  (\url{https://classroom.google.com}).
  \item{GNU Bash Manual} --- \url{https://www.gnu.org/software/bash/manual/}
  \item{Coreutils documentation} --- \url{https://www.gnu.org/software/coreutils/}
  \item{Docker documentation} --- \url{https://docs.docker.com/}
  \item{Pro Git (online)} --- \url{https://git-scm.com/book}
  % TODO: пополнять.
\end{description}

% Приложение: Дорожная карта курса
\chapter{Дорожная карта курса}
\label{app:roadmap}

Дорожная карта показывает, какие темы методички связаны с лекциями и
лабораторными работами, и фиксирует статус наполнения. Заполняется по мере
прохождения курса.

\begin{longtable}{@{}p{0.45\linewidth}p{0.35\linewidth}l@{}}
\toprule
Тема методички & Занятие (лекция / лаб.) & Статус \\
\midrule
\endhead

Доступ к окружению \ref{ch:02-linux-access}         & Лаб.\ 1 (сем.\ 1)      & готово \\
Оболочка и файловая система \ref{ch:03-shell}       & Лаб.\ 1--2, лек.\ 2   & готово \\
Права доступа и пользователи \ref{ch:05-permissions} & Лаб.\ 2 (сем.\ 1)   & готово \\
Справка, конвейеры, пакеты \ref{ch:07-text-pipelines} & Лаб.\ 2--3, лек.\ 2 & готово \\
Редакторы и сессии (vim/tmux/mc) \ref{ch:09-vim}    & Лаб.\ 2, лек.\ 2      & готово \\
Процессы и ресурсы \ref{ch:12-processes}            & Лек.\ 1 (сем.\ 1)    & готово \\
C, сборка, тесты, память, отладка \ref{ch:13-c-compilation} & Лек.\ 1, лаб.\ (C) & готово \\
Docker и изоляция \ref{ch:19-docker-basics}         & Лек.\ 1, лаб.\ (Docker) & готово \\
Git \ref{ch:25-git}                                  & Лаб.\ 4 (сем.\ 1)    & готово \\
Основы безопасности \ref{ch:24-security}            & Лек.\ (сем.\ 2)      & задел \\
Пайплайны / CI \ref{ch:26-ci}                        & Лек.\ 2 (сем.\ 2)    & задел \\
Скриптинг / IDE \ref{ch:27-scripting-ide}           & ---                  & задел \\
\bottomrule
\end{longtable}

% Легенда статусов: «готово» = раздел наполнен; «будет дополнено» = заглушка,
% материал появится с новыми семинарами/лекциями.


\end{document}
